\documentclass[11pt,a4paper]{article}
\usepackage[utf8]{inputenc}
\usepackage[T1]{fontenc}
\usepackage{textcomp}
\usepackage[margin=2.5cm]{geometry}
\usepackage{booktabs}
\usepackage{amsmath}
\usepackage{amssymb}
\usepackage{siunitx}
\usepackage{multirow}
\usepackage{float}
\usepackage{longtable}
\usepackage{array}
\usepackage[autostyle]{csquotes}
\usepackage{xcolor}
\usepackage{listings}
\usepackage[hidelinks]{hyperref}
\usepackage{graphicx}

\sisetup{detect-weight=true}

\lstdefinelanguage{toml}{morecomment=[l]{\#}, morestring=[b]"}
\lstdefinestyle{base}{basicstyle=\ttfamily\scriptsize, numbers=left,
  numberstyle=\tiny\color{gray}, numbersep=6pt, breaklines=true,
  breakatwhitespace=false, columns=fullflexible, keepspaces=true,
  showstringspaces=false, frame=single, framerule=0.3pt, rulecolor=\color{gray!60},
  xleftmargin=2.2em, framexleftmargin=0.3em, tabsize=4, upquote=true,
  captionpos=t, aboveskip=6pt, belowskip=8pt}
\lstdefinestyle{py}{style=base, language=Python, keywordstyle=\color{blue!60!black},
  commentstyle=\color{gray!75!black}, stringstyle=}
\lstdefinestyle{sh}{style=base, language=bash, keywordstyle=\color{blue!60!black},
  commentstyle=\color{gray!75!black}, stringstyle=}
\lstdefinestyle{toml}{style=base, language=toml, commentstyle=\color{gray!75!black},
  stringstyle=}
\renewcommand{\lstlistingname}{Code}
\newcommand{\cf}[1]{\texttt{#1}}

\title{\bfseries Marocco FL thesis --- SAGE-smart: complete specification with code}
\author{}
\date{}

\begin{document}
\maketitle
\vspace{-3.2em}
\begin{center}
Client-Driven Local Work Assignment in Battery-Constrained Federated Learning
\end{center}
\vspace{1em}

\tableofcontents
\newpage

%======================================================================
\section{Scope and conventions}
\label{sec:scope}

This document describes every rule, parameter and computation of SAGE-smart as
it is implemented in the simulator, together with the code that implements it.
SAGE-smart is the protocol that the reports up to week 09/11 called
SAGE-smart 2. Since 16/11/2026 it is the only SAGE-smart in the code: the
previous versions were removed. The labels \cf{sage\_smart\_lin},
\cf{sage\_smart\_peuk} and \cf{sage\_smart\_nm} run the same algorithm. They
differ only in the battery model of the simulated phones: Ideal (linear),
Peukert and Datasheet.

The code excerpts are copied automatically from the source files, with their
real line numbers. Line numbers refer to the code of 16/11/2026. Paths are
relative to the root of the repository. Comments inside the code are in
Italian, as in the sources.

The implementation has three parts.
\begin{itemize}
  \item Server: the strategy \cf{PhysicalSAGESmart}
        (\cf{tesifl/strategies/sage\_smart.py}). It inherits the selection
        skeleton and the label divergence from \cf{PhysicalSAGEAblation}
        (\cf{strategies/sage\_ablation.py}), and the round orchestration and
        the physical accounting from \cf{PhysicalFedAvg}
        (\cf{strategies/fedavg.py}), which extends the Flower \cf{FedAvg}.
  \item Client: the Flower \cf{ClientApp} (\cf{tesiFL/client\_app.py}), the
        decision policy (\cf{tesiFL/client\_policy.py}) and the view of the
        phone on its own state (\cf{device/fuel\_gauge.py}).
  \item Physical world: \cf{device/world\_state.py} and the models it uses
        (compute, communication, battery, workload, recharge). It simulates
        the phones and lives in the server process. It is not part of the
        algorithm: the algorithm reads it only through the channels listed in
        Section~\ref{sec:info}. The server estimates and the client
        predictions use the same model functions as the world
        (Section~\ref{sec:physics}).
\end{itemize}

\subsection*{Notation}

\begin{longtable}{p{0.20\textwidth}p{0.74\textwidth}}
\toprule
Symbol & Meaning \\
\midrule
\endhead
$\mathcal{N}$, $N$ & set of phones, $N = |\mathcal{N}| = 100$; $i$ indexes a phone \\
$r$ & round, $r = 1, \dots, 250$ \\
$\mathcal{A}_r$, $m_r$ & phones alive at round $r$ (battery not depleted), $m_r = |\mathcal{A}_r|$ \\
$S_r$ & phones that receive the model at round $r$, $S_r \subseteq \mathcal{A}_r$ \\
$n_i$ & training samples of phone $i$ \\
$|\mathcal{Y}|$ & number of classes, 100 (CIFAR-100) \\
$E$, $B$ & local epochs and batch size; $E_0 = 5$, $B_0 = 128$ of the baselines \\
$E_{\mathrm{lo}}$, $E_{\mathrm{cap}}$, $E_{\mathrm{hi},i}$ & minimum epochs (2), maximum epochs (10), maximum epochs of phone $i$ in this round \\
$\mathcal{B}$ & batch sizes of the envelope, $\{16, 32, 64, 128\}$ \\
$T$ & round deadline \\
$s_i$, $\hat{s}_i$ & true state of charge (SoC) of phone $i$, SoC known to the server \\
$s_{\min}$ & battery reserve, 0.20 \\
$w_i$, $\hat{w}_i$ & workload of the other apps (fraction of compute capacity), value known to the server \\
$\bar{w}$, $w_{\max}$ & population mean workload (0.118), heaviest workload state (0.30) \\
$\mu_i$ & compute throughput of the tier, MAC/s \\
$P^{\mathrm{tr}}_i$, $P^{\mathrm{idle}}_0$ & training power of the tier, idle power (0.1 W) \\
$M$, $F$, $o$ & forward MACs per sample, backward factor (3), overhead per batch (0.057 s) \\
$Q_i$, $V_0$, $\eta$ & nominal capacity (Ah), nominal voltage (3.7 V), DC-DC efficiency (1.0) \\
$\lambda$, $\gamma$ & recharge rate of the users (0.126 per hour), charging speed (1 SoC per hour, 1C) \\
$a$, $b$, $c$ & score weights of SoC, divergence and staleness \\
$D_i$, $\theta_i$, $\ell_i$ & label divergence, staleness, last round in which $i$ was selected or refused \\
$\Theta$ & staleness scale, 25 rounds \\
$\sigma_i$, $\phi_i$ & ranking score, fresh score \\
$\kappa$, $W$ & shortlist factor (1.5), maximum number of proposal waves (2) \\
$p_{\mathrm{nr}}$, $t_{\mathrm{to}}$, $t_{\mathrm{ctrl}}$ & non-response probability (0.10), timeout (1 s), control-message latency (0.335 s) \\
$k_i$, $\alpha$ & energy calibration factor of phone $i$, its smoothing weight (0.3) \\
$\tau$, $\delta$ & relaxation time of the workload forecast (450 s), forecast step ($\le 300$ s) \\
$L_i$, $u_i$, $h_i$ & probe loss, statistical utility, battery headroom of phone $i$ \\
$\nu_i$ & SGD steps executed by phone $i$ \\
\bottomrule
\end{longtable}

\newpage
%======================================================================
\section{Scenario and parameters}
\label{sec:params}

\subsection{Scenario}

The campaign runs $N = 100$ phones for 250 rounds. The model is ResNet-18 in its
CIFAR variant (3$\times$3 initial convolution, no max-pool) on CIFAR-100: $M =
555{,}468{,}800$ forward MACs per sample and 11{,}220{,}132 parameters, that is
44.88 MB in float32. The training set is split among the phones with the
Dirichlet partitioner of \cf{flwr-datasets} (\cf{alpha} $= \beta \in \{0.1,
0.5, 1.0\}$, same seed as the run). Each partition is split again 80/20 into
local training and local test; $n_i$ is the training part, 400 samples on
average. The client trains with random crop (padding 4) and horizontal flip,
SGD with learning rate 0.01 and momentum 0.9, cross-entropy loss. The global
model is evaluated on the full CIFAR-100 test set after every round.

\begin{lstlisting}[style=py,firstnumber=98,caption={\texttt{tesifl/tesiFL/data/partition.py}, lines 98--122}]
def partition_sizes(num_partitions: int, beta: float, seed: int,
                    dataset: str = "cifar10") -> dict[int, int]:
    """Numero di campioni di TRAINING di ogni partizione, senza caricare le
    immagini in memoria piu' del necessario.

    Riproduce esattamente lo split di load_data (stesso partitioner, stesso
    train_test_split con test_size=0.2 e stesso seed), quindi il valore
    coincide con `num-examples` che il client riporta a fine training.

    Serve alle strategie che, come ESCS, hanno bisogno della latenza di
    training di OGNI client prima del primo round (nel paper e' un dato di
    profilo, `total_train_latency`).
    """
    hf_name, label_col, _ = _spec(dataset)
    partitioner = DirichletPartitioner(
        num_partitions=num_partitions, partition_by=label_col, alpha=beta, seed=seed,
    )
    fds_local = FederatedDataset(
        dataset=hf_name, partitioners={"train": partitioner},
    )
    sizes: dict[int, int] = {}
    for pid in range(num_partitions):
        part = fds_local.load_partition(pid)
        sizes[pid] = len(part.train_test_split(test_size=0.2, seed=seed)["train"])
    return sizes
\end{lstlisting}

\subsection{Where the parameters come from}

The only file to edit is \cf{experiment.toml}. \cf{run.sh} reads it and passes
the values to Flower as run configuration (\cf{--run-config}); the defaults in
\cf{pyproject.toml} are always overridden. \cf{\_shard.sh} adds the score
weights of the current $\beta$ and the battery of the world, taken from the
suffix of the label. \cf{server\_app.py} builds the strategy.

\begin{lstlisting}[style=toml,firstnumber=290,caption={\texttt{tesifl/experiment.toml}, lines 290--350}]
# Ablazione di SAGE (sage_peuk, sage_nm) e SAGE-smart: pesi (a,b) per beta,
# dalla sweep in results_findbest.
# [B] SAGE-smart usa gli stessi dell'ablazione: il confronto fra i due isola
# l'assegnazione dinamica di (E,B), non la taratura.
sage_ab  = { "0.1" = [0.4, 0.6], "0.5" = [0.6, 0.4], "1.0" = [0.9, 0.1] }

# SAGE-smart: il server seleziona (a*SoC + b*D + c*anzianita') e propone un
# envelope di (E, B); il CLIENT sceglie dentro l'envelope con il proprio stato.
# Etichette sage_smart_lin / _peuk / _nm: stesso algoritmo, cambia solo la
# batteria dei device.
smart_stale_weight   = 0.3    # peso (c) dell'anzianita' nel punteggio
smart_stale_max      = 25     # round dopo i quali l'anzianita' e' al massimo
smart_soc_min        = 0.20   # riserva che il client non intacca (vincolo 2e di SAGE)
smart_epochs_min     = 2      # E_lo: minimo di epoche dell'envelope
smart_deadline_mult  = 1.0    # deadline = mult * mediana dei tempi nominali (E, B di [algo])
smart_probe_samples  = 256    # campioni locali per stimare la loss prima di decidere
# Negoziazione prima del training (documento "Client Server Protocol" del
# relatore, opzione 3 con l'opzione 2): il server ordina i client con la
# previsione del SoC, propone l'envelope a una rosa di 1.5 volte i posti e
# sceglie sullo stato appena riportato; se le accettazioni non bastano, una
# seconda ondata (chi rifiuta lo fa solo per energia).
smart_max_waves      = 2      # ondate di proposte per round (1 = nessun rimpiazzo)
# B di SAGE-smart (le baseline usano batch_size = 128). L'envelope propone le
# potenze di 2 fra smart_batch_min e smart_batch_max; il client sceglie il
# B "piu' economico per passo di SGD": il piu' piccolo, riducendo prima E, e
# sale al successivo solo se a quel B non ci sta nemmeno E_lo (batteria o
# deadline). Vedi tesiFL/client_policy.py.
# [B] B_min = 16 dallo sweep (results_sweepB*, SAGE-smart nm, 3 beta, B fisso):
#   B     8      16     32     64     128    256
#   acc   0.702  0.667  0.632  0.608  0.586  0.531
#   SoC   0.500  0.543  0.575  0.574  0.578  0.576
# l'accuracy cresce sempre al calare di B; il SoC e' piatto da 32 in su e
# cala sotto. 16 e' il punto piu' vicino all'ideale (acc e SoC massimi).
# [B] TETTO A 64, non 256 (14/10, quando SAGE-smart usava FedNova: spento dal
# 30/10, tolto dal codice il 16/11). SAGE-smart aggregava con FedNova, che amplifica
# l'update di chi fa pochi passi di tau_eff/tau_i. Col range 16-256 un client
# fuori deadline ripiegava su B=256 e faceva ~18 passi contro ~290: spread
# 15-16 (soglia ~6-8) e modello collassato nel fastTest del 14/10. A 64 lo
# stesso client fa ~70 passi. Le partizioni grandi salgono a 32/64 e fanno
# comunque un numero di passi simile ai client a 16 (3000/32 ~ 1666/16).
smart_batch_min      = 16
smart_batch_max      = 128    # 64 per ResNet-20; 128 richiesto per ResNet-18 (SCENARIO)

# BROADCAST DELLO STATO a inizio round (opzione 1 del "Client Server
# Protocol" del relatore, call del 16/10) per le strategie che selezionano
# sullo stato: SAGE, ESCS, SAGE-smart (non FedAvg/FedProx, che scelgono a
# caso). Nel simulatore non cambia l'informazione (SAGE/ESCS leggono il SoC
# del mondo, SAGE-smart la sua stima, errore < 1 punto di SoC): serve a
# CONTARNE IL COSTO, perche' dichiariamo di avere un SoC fidato a inizio
# round. Chi risponde paga uno scambio di controllo (sessione 4G intera,
# ~2.1 J, se non viene selezionato; il solo payload se viene selezionato);
# chi non risponde e' escluso dal round; il round si allunga della latenza
# o del timeout. Stima: ~3 Wh di broadcast + ~1 Wh di attesa su 250 round.
# [B] 16/11: SAGE-smart NON usa il broadcast (lo sostituisce la negoziazione:
# le risposte alle proposte danno raggiungibilita' e stato; broadcast_nonresp
# e broadcast_timeout_s valgono per le proposte). Lo usa solo l'ablazione
# senza negoziazione (noneg).
broadcast            = true
broadcast_nonresp    = 0.10   # client che non rispondono (indicazione del relatore)
broadcast_timeout_s  = 1.0    # [ASSUNZIONE] attesa delle risposte mancanti (~3x promozione + RTT)
broadcast_top_m      = 0      # [B] 16/11: non lo usa piu' nessuno (era solo di SAGE-smart)
\end{lstlisting}

\begin{lstlisting}[style=sh,firstnumber=73,caption={\texttt{tesifl/run.sh}, lines 73--95}]
        "ALGOCFG": (f'model-name="{mo["name"]}" dataset="{mo["dataset"]}" '
                f"augment={str(mo.get('augment', False)).lower()} "
                f"local-epochs={al['local_epochs']} batch-size={al['batch_size']} "
                f"learning-rate={al['learning_rate']} "
                f"fraction-evaluate={al['fraction_evaluate']} "
                f"proximal-mu={al['proximal_mu']} "
                f"smart-stale-weight={al['smart_stale_weight']} "
                f"smart-stale-max={al['smart_stale_max']} "
                f"smart-soc-min={al['smart_soc_min']} "
                f"smart-epochs-min={al['smart_epochs_min']} "
                f"smart-deadline-mult={al['smart_deadline_mult']} "
                f"smart-probe-samples={al['smart_probe_samples']} "
                f"smart-max-waves={al['smart_max_waves']} "
                f"smart-batch-min={al['smart_batch_min']} "
                f"smart-batch-max={al['smart_batch_max']} "
                f"broadcast={str(al['broadcast']).lower()} "
                f"broadcast-nonresp={al['broadcast_nonresp']} "
                f"broadcast-timeout-s={al['broadcast_timeout_s']} "
                f"broadcast-top-m={al['broadcast_top_m']} "
                f"escs-min-battery={al['escs_min_battery']} "
                f"escs-min-nq={al['escs_min_nq']} "
                f"escs-first-round-all={str(al['escs_first_round_all']).lower()} "
                f"escs-cap-probabilistic={str(al['escs_cap_probabilistic']).lower()}"),
\end{lstlisting}

\begin{lstlisting}[style=sh,firstnumber=23,caption={\texttt{tesifl/\_shard.sh}, lines 23--32}]
# pesi di SAGE per beta, letti da experiment.toml (o da EXP_TOML, vedi run.sh)
read -r SA SB SC QA QB <<< "$(python3 - "$BETA" "${EXP_TOML:-experiment.toml}" <<'PY'
import sys, tomllib
b = sys.argv[1]
c = tomllib.load(open(sys.argv[2], "rb"))["algo"]
abc = c["sage_abc"].get(b, [0.5, 0.2, 0.3])
ab = c["sage_ab"].get(b, [0.5, 0.5])
print(*abc, *ab)
PY
)"
\end{lstlisting}

\begin{lstlisting}[style=sh,firstnumber=77,caption={\texttt{tesifl/\_shard.sh}, lines 77--113}]
      # [B] 16/11: sage_smart e' l'ex SAGE-smart 2 (previsione del SoC,
      # negoziazione con rosa e scelta fresca, niente broadcast, lavoro
      # parziale), l'unico SAGE-smart: fra _lin, _peuk e _nm cambia solo la
      # batteria dei device. Le altre etichette SPENGONO componenti
      # (smart-ablation in server_app): nostale, noest, noneg, e5, b128,
      # nopartial.
      sage_smart) algo="sage_smart"; extra="sage-a=${QA} sage-b=${QB}" ;;
      # ablazione leave-one-out: un componente spento, il nome e' il token
      sage_smart_nostale|sage_smart_noest|sage_smart_noneg|sage_smart_e5|sage_smart_b128|sage_smart_nopartial)
                  algo="sage_smart"
                  extra="sage-a=${QA} sage-b=${QB} smart-ablation='${base#sage_smart_}'" ;;
      # ablazione CUMULATA (ablation_cum.sh), nomi delle run gia' fatte.
      # Ordine diritto: 0 cum0, 1 cum1 (+ anzianita'), 2 cum2p (+ E adattivo
      # col lavoro parziale), 3 pcum4 (+ B adattivo), 4 sage_smart2_nofc
      # (+ negoziazione), 5 sage_smart (+ previsione del SoC).
      # Ordine inverso: 0 cum0, 1 cum1, 2 sage_smart2_rev2 (+ negoziazione e
      # previsione, E = 5 e B = 128 fissi), 3 sage_smart2_rev3 (+ E
      # adattivo), 4 sage_smart (+ B adattivo).
      # [B] 16/11: cum2p e pcum4 erano girate col carico MASSIMO
      # all'iscrizione (tolto): col codice di oggi non tornano identiche.
      # Le altre sono uguali alle run fatte fino al 12/11.
      sage_smart_cum0) algo="sage_smart"
                  extra="sage-a=${QA} sage-b=${QB} smart-ablation='nostale+noest+noneg+e5+b128+nopartial'" ;;
      sage_smart_cum1) algo="sage_smart"
                  extra="sage-a=${QA} sage-b=${QB} smart-ablation='noest+noneg+e5+b128+nopartial'" ;;
      sage_smart_cum2p) algo="sage_smart"
                  extra="sage-a=${QA} sage-b=${QB} smart-ablation='noest+noneg+b128'" ;;
      sage_smart_pcum4) algo="sage_smart"
                  extra="sage-a=${QA} sage-b=${QB} smart-ablation='noest+noneg'" ;;
      sage_smart2_nofc) algo="sage_smart"
                  extra="sage-a=${QA} sage-b=${QB} smart-ablation='noest'" ;;
      sage_smart2_rev2) algo="sage_smart"
                  extra="sage-a=${QA} sage-b=${QB} smart-ablation='e5+b128+nopartial'" ;;
      sage_smart2_rev3) algo="sage_smart"
                  extra="sage-a=${QA} sage-b=${QB} smart-ablation='b128'" ;;
      # nome di sage_smart nelle cumulate fatte fino al 12/11 (stessa run)
      sage_smart2_socfc) algo="sage_smart"; extra="sage-a=${QA} sage-b=${QB}" ;;
\end{lstlisting}

The suffix \cf{\_lin}, \cf{\_peuk} or \cf{\_nm} of the label becomes
\cf{world-battery}, and a suffix \cf{\_x<m>} becomes \cf{workload-scale} (1 in
the campaign). The strategy is built as follows; the values not passed here
are the defaults of the constructor (Section~\ref{sec:init}).

\begin{lstlisting}[style=py,firstnumber=174,caption={\texttt{tesifl/tesiFL/server\_app.py}, lines 174--249}]
    elif algorithm == "sage_smart":
        # [B] la taglia dei dataset e' un dato di ISCRIZIONE (Fase 0 della
        # proposta): serve al server per la deadline e per l'envelope prima
        # che il client abbia mai risposto. Stesso partitioner di ESCS.
        sizes = partition_sizes(
            num_partitions=n_clients,
            beta=float(context.run_config["beta"]),
            seed=seed,
            dataset=_DATASET,
        )
        smart = dict(
            stale_weight=float(context.run_config["smart-stale-weight"]),
            epochs_min=int(context.run_config["smart-epochs-min"]),
            batch_min=int(context.run_config["smart-batch-min"]),
            batch_max=int(context.run_config["smart-batch-max"]),
        )
        # [B] 16/11: SAGE-smart e' l'ex SAGE-smart 2 (vedi la docstring di
        # PhysicalSAGESmart): previsione del SoC, negoziazione con rosa e
        # scelta fresca, niente broadcast, lavoro parziale sono i default.
        # smart-ablation SPEGNE componenti, per le ablazioni (piu' componenti
        # uniti da "+", es. "noest+noneg"). Un solo parametro (e non i singoli
        # valori) perche' --run-config non ammette chiavi ripetute e quelle
        # singole sono gia' in ALGOCFG.
        ablation = str(context.run_config["smart-ablation"])
        ablations = {
            "none":      {},
            "nostale":   {"stale_weight": 0.0},                       # niente c
            "noest":     {"forecast": False},                         # SoC dell'ultimo report
            "noneg":     {"negotiation": False},                      # niente proposte
            "e5":        {"epochs_min": epochs, "epochs_cap_mult": 1},  # E = 5 fisso
            "b128":      {"batch_min": batch_size, "batch_max": batch_size},  # B = 128
            "nopartial": {"partial_work": False},                     # niente lavoro parziale
        }
        parts = [p for p in ablation.split("+") if p and p != "none"]
        unknown = [p for p in parts if p not in ablations]
        if unknown:
            raise ValueError(f"smart-ablation sconosciuta: {'+'.join(unknown)!r} "
                             f"(usa {', '.join(ablations)}, uniti da '+')")
        for p in parts:
            smart.update(ablations[p])
        if smart.get("negotiation", True):
            # chi e' raggiungibile e in che stato lo dicono le risposte alle
            # proposte: niente broadcast, la frazione broadcast-nonresp non
            # risponde alle proposte
            bcast = dict(bcast, broadcast=False)
            print("[strategy] SAGE-smart: niente broadcast dello stato, "
                  f"{bcast['broadcast_nonresp']:.0%} di non risposte sulle proposte")
        elif bcast["broadcast"]:
            # ablazione noneg: senza proposte torna il broadcast dello stato
            print(f"[strategy] SAGE-smart senza negoziazione: broadcast dello stato, "
                  f"{bcast['broadcast_nonresp']:.0%} non risponde, timeout "
                  f"{bcast['broadcast_timeout_s']} s")
        strategy = PhysicalSAGESmart(
            sage_a=float(context.run_config["sage-a"]),
            sage_b=float(context.run_config["sage-b"]),
            stale_max=int(context.run_config["smart-stale-max"]),
            partition_sizes=sizes,
            soc_min=float(context.run_config["smart-soc-min"]),
            deadline_mult=float(context.run_config["smart-deadline-mult"]),
            probe_samples=int(context.run_config["smart-probe-samples"]),
            max_waves=int(context.run_config["smart-max-waves"]),
            **smart, **common, **bcast,
        )
        print(
            f"[strategy] PhysicalSAGESmart ({strategy.variant}) "
            f"ablazione {ablation} | "
            f"(a={context.run_config['sage-a']}, "
            f"b={context.run_config['sage-b']}, "
            f"c={strategy.stale_w}) | "
            f"(E,B) decisi dal client | "
            f"soc_min={strategy.soc_min} E {strategy.epochs_min}-{strategy.epochs_cap} "
            f"deadline x{strategy.deadline_mult} probe={strategy.probe_samples} "
            + (f"rosa di {strategy.shortlist:g} k, max {strategy.max_waves} ondate | "
               if strategy.negotiation else "")
            + f"B in {strategy.batch_choices}"
        )
\end{lstlisting}

\subsection{Parameters of SAGE-smart}

\begin{longtable}{p{0.24\textwidth}p{0.08\textwidth}p{0.62\textwidth}}
\toprule
Parameter (code) & Value & Meaning and source \\
\midrule
\endhead
\cf{sage\_a}, \cf{sage\_b} ($a$, $b$) & see right & SoC and divergence weights, $a + b = 1$: (0.4, 0.6) at $\beta = 0.1$, (0.6, 0.4) at 0.5, (0.9, 0.1) at 1.0. \cf{sage\_ab}, \cf{experiment.toml} line 294; same weights as the SAGE ablation \\
\cf{stale\_weight} ($c$) & 0.3 & staleness weight; \cf{smart\_stale\_weight}, line 300 \\
\cf{stale\_max} ($\Theta$) & 25 & rounds after which the staleness saturates at 1; line 301 \\
\cf{soc\_min} ($s_{\min}$) & 0.20 & reserve protected by the client; line 302 \\
\cf{epochs\_min} ($E_{\mathrm{lo}}$) & 2 & minimum epochs of the envelope; line 303 \\
\cf{epochs} ($E_0$) & 5 & epochs of the baselines; reference work for the deadline; \cf{local\_epochs}, line 282 \\
\cf{epochs\_cap\_mult} & 2 & $E_{\mathrm{cap}} = 2 E_0 = 10$; constructor default \\
\cf{batch\_size} ($B_0$) & 128 & batch of the baselines; reference work for the deadline; line 283 \\
\cf{batch\_min}, \cf{batch\_max} & 16, 128 & $\mathcal{B}$ = powers of 2 between them, $\{16, 32, 64, 128\}$; lines 330--331 \\
\cf{deadline\_mult} & 1.0 & $T$ = this factor times the median nominal time; line 304 \\
\cf{probe\_samples} & 256 & samples on which the client measures the loss; line 305 \\
\cf{max\_waves} ($W$) & 2 & proposal waves per round; line 311 \\
\cf{shortlist} ($\kappa$) & 1.5 & proposals to $\lceil \kappa \cdot \text{slots} \rceil$ phones; constructor default \\
\cf{load\_tau\_s} ($\tau$) & 450 s & relaxation time of the workload forecast; constructor default \\
\cf{calib\_alpha} ($\alpha$) & 0.3 & weight of the moving average of $k_i$; constructor default \\
\cf{query\_timeout\_s} & 600 s & maximum wait of the Flower \cf{send\_and\_receive} for the replies to the proposals (simulator safeguard); constructor default \\
\cf{sage\_k} ($k$) & 2 & slots at round 1; default of \cf{PhysicalSAGEAblation} \\
\cf{max\_fraction} & 0.2 & slots are at most 20\% of the alive phones; default of \cf{PhysicalSAGEAblation} \\
\cf{broadcast\_nonresp} ($p_{\mathrm{nr}}$) & 0.10 & probability that a contacted phone does not answer; line 348 \\
\cf{broadcast\_timeout\_s} ($t_{\mathrm{to}}$) & 1.0 s & wait for missing answers; line 349 \\
\cf{broadcast} & off & on in \cf{experiment.toml} for SAGE and ESCS; \cf{server\_app} switches it off for SAGE-smart, which uses the negotiation instead \\
\cf{forecast}, \cf{negotiation}, \cf{partial\_work} & on & components that the ablations switch off (Section~\ref{sec:ablation}) \\
\cf{\_ESTIMATE\_STEP\_S} & 300 s & maximum integration step of the forecast; \cf{sage\_smart.py} line 27 \\
\cf{learning\_rate} & 0.01 & SGD of the client (momentum 0.9); line 284 \\
\bottomrule
\end{longtable}

\subsection{Parameters of the world used by the computations of SAGE-smart}

\begin{longtable}{p{0.30\textwidth}p{0.22\textwidth}p{0.40\textwidth}}
\toprule
Parameter & Value & Source (\cf{experiment.toml}) \\
\midrule
\endhead
tiers: capacity $Q$, throughput $\mu$, training power $P^{\mathrm{tr}}$ & low 4820 mAh, 13.9 GMAC/s, 2.99 W; mid 4780 mAh, 26.6 GMAC/s, 4.00 W; high 5000 mAh, 30.5 GMAC/s, 6.15 W & \cf{[device].tiers}, line 125; one third of the phones each, permuted with the seed \\
idle power $P^{\mathrm{idle}}_0$ & 0.1 W & line 136 \\
initial SoC & normal (0.75, 0.15) truncated to [0.30, 1.00] & lines 150--152 \\
$V_0$, rated C-rate, $\eta$, Peukert $n$ & 3.7 V, 0.2C, 1.0, 1.05 & \cf{[battery]}, lines 156--167 \\
$F$, $o$ & 3.0, 0.057 s & \cf{[compute]}, lines 173, 185 \\
bytes per parameter & 4 & line 193 \\
4G energy per MB, down / up & 0.98 / 2.96 J/MB & lines 207--208 \\
4G bandwidth, down / up & 25.1 / 10.0 Mbit/s & lines 209--210 \\
radio session: promotion & 1.2107 W for 0.265 s & lines 223--224 \\
radio session: tail & 0.178 W for 10.2 s & lines 225--226 \\
RTT, control payload & 69.5 ms, 1 KB & lines 227--228 \\
workload states & 0, 0.10, 0.30 & \cf{[workload]}, line 242 \\
workload chains (low, medium, high) & three Markov chains, stationary means 0.030, 0.119, 0.206 & lines 247--249 \\
recharge rate $\lambda$, charging speed $\gamma$ & 0.126 per hour, 1C & \cf{[recharge]}, lines 270, 275 \\
\bottomrule
\end{longtable}

\newpage
%======================================================================
\section{Physical models used inside SAGE-smart}
\label{sec:physics}

The server uses these functions to compute the deadline, the envelope and the
forecast; the client uses them to predict the time and the SoC at the end of
its work. They are the same functions with which the world drains the
batteries, so the client prediction of its own final SoC is exact (CSV column
\cf{smart\_soc\_err} $\approx 0$) and the server forecast is wrong only because
of the recharges and the workload it cannot see.

\subsection{Compute time}

For $E$ epochs (possibly fractional) on $n$ samples with batch $B$:
\begin{equation}
  t^{\mathrm{tr}}(E, n, B) = \frac{F \, M \, E \, n}{\mu^{\mathrm{eff}}}
                           + E \cdot \max\!\bigl(1, \lfloor n / B \rfloor\bigr) \cdot o ,
  \qquad
  \mu^{\mathrm{eff}} = \mu \cdot \max(10^{-3}, 1 - w) .
  \label{eq:ttrain}
\end{equation}
The loss probe on $n_{\mathrm{p}}$ samples is a forward pass only:
$t^{\mathrm{probe}} = M n_{\mathrm{p}} / \mu^{\mathrm{eff}}$. The workload $w$
slows the training down; the training power does not change. With $n = 400$
and $w = \bar{w}$ one epoch lasts 55.8 s (low), 29.8 s (mid) and 26.2 s
(high) at $B = 16$, and 54.5, 28.6 and 25.0 s at $B = 128$. The probe on 256
samples lasts 11.6, 6.1 and 5.3 s.

\begin{lstlisting}[style=py,firstnumber=11,caption={\texttt{tesifl/device/compute\_model.py}, lines 11--41}]
def training_macs(epochs: int, n_samples: int) -> float:
    #MACs totali di un round locale: fwd*factor x campioni x epoche
    return TRAIN_MACS_FACTOR * MODEL_MACS_FWD_PER_SAMPLE * epochs * n_samples


def training_time_s(
    profile: DeviceProfile, epochs: int, n_samples: int, batch_size: int
) -> float:
    """Durata di un round: calcolo + overhead per batch.

    [B] l'overhead fisso per batch e' l'unico modo in cui B entra nel costo:
    la potenza di training e' costante per tier (DeviceProfile.train_power_w),
    quindi l'energia di un'epoca, P * t, CALA al crescere di B. Non c'e' un
    batch energeticamente ottimo interno: la scelta di B e' una questione di
    accuracy (passi di SGD per epoca), che il modello fisico non vede.
    """
    n_batches = max(1, n_samples // batch_size)
    compute_s = training_macs(epochs, n_samples) / profile.macs_per_s
    overhead_s = epochs * n_batches * BATCH_OVERHEAD_S
    return compute_s + overhead_s


def probe_time_s(profile: DeviceProfile, n_samples: int) -> float:
    """Solo forward su n_samples: la stima della loss locale prima del training.

    [B] niente fattore 3 (non c'e' backward) e niente overhead per batch: e'
    una singola passata in valutazione, trascurabile rispetto al resto.
    """
    if n_samples <= 0:
        return 0.0
    return MODEL_MACS_FWD_PER_SAMPLE * n_samples / profile.macs_per_s
\end{lstlisting}

\subsection{Power}

During training a phone draws the power of its tier, $P^{\mathrm{tr}}_i$,
whatever $B$ is. Between rounds it draws
\begin{equation}
  P^{\mathrm{idle}}(w) = P^{\mathrm{idle}}_0 + w \, (P^{\mathrm{tr}}_i - P^{\mathrm{idle}}_0) ,
  \label{eq:pidle}
\end{equation}
0.441, 0.561 and 0.815 W at $w = \bar{w}$ for the three tiers.

\begin{lstlisting}[style=py,firstnumber=37,caption={\texttt{tesifl/device/world\_state.py}, lines 37--63}]
def effective_profile(p: DeviceProfile, w: float) -> DeviceProfile:
    """Profilo col throughput ridotto dal workload w."""
    if w <= 0.0:
        return p
    return DeviceProfile(
        cid=p.cid, battery_capacity_mah=p.battery_capacity_mah,
        macs_per_s=p.macs_per_s * max(1e-3, 1.0 - w),
        peukert_n=p.peukert_n, tier=p.tier, train_power_w=p.train_power_w,
    )


def training_power_w(p: DeviceProfile) -> float:
    """Potenza del device durante il training: costante, quella del suo tier.

    [B] prima P(B) = 3 W * (B/128)^0.3, uguale per tutti i tier e con un
    esponente senza fonte. Ora la potenza e' quella MISURATA del telefono
    del tier (Mandheling, vedi experiment.toml) e non dipende da B.
    """
    return p.train_power_w


def idle_power_w(p: DeviceProfile, w: float) -> float:
    """P_idle + w * (P_train - P_idle): interpola fra riposo e pieno carico
    DEL TIER. w e' la frazione della capacita' di quel telefono occupata da
    altre app: a parita' di w un telefono piu' potente fa piu' lavoro e
    consuma piu' watt."""
    return IDLE_POWER_W + w * (p.train_power_w - IDLE_POWER_W)
\end{lstlisting}

\subsection{Communication}

The model is $S = 11{,}220{,}132 \times 4\,\text{B} = 44.88$ MB, the same size in
both directions. Every isolated transfer opens a 4G radio session (promotion and
tail):
\begin{align}
  E^{\mathrm{sess}} &= 1.2107 \cdot 0.265 + 0.178 \cdot 10.2 = 2.136~\text{J},
  & t^{\mathrm{sess}} &= 0.265 + 10.2 = 10.465~\text{s}, \\
  E^{\mathrm{comm}} &= 0.98\,S + 2.96\,S + 2 E^{\mathrm{sess}} = 181.10~\text{J},
  & t^{\mathrm{comm}} &= \frac{S}{25.1/8} + \frac{S}{10/8} + 2 t^{\mathrm{sess}} = 71.14~\text{s}.
\end{align}
The battery models work with power and time, so the model exchange is applied
as the average power $P^{\mathrm{comm}} = E^{\mathrm{comm}} / t^{\mathrm{comm}}
= 2.546$ W for $t^{\mathrm{comm}}$. A phone that downloads the model and then
withdraws pays only the download: 46.12 J in 24.77 s. A control message
(proposal and answer, 1 KB, half in each direction) costs $1.97 \cdot 10^{-3}$ J
when it shares the radio session of the model download that follows (selected
phone), and 2.138 J when it is isolated. It delays the round by
$t_{\mathrm{ctrl}} = 0.265 + 0.0695 + 0.0006 = 0.335$ s (promotion, RTT,
transmission); the tail does not delay the round.

\begin{lstlisting}[style=py,firstnumber=44,caption={\texttt{tesifl/device/communication\_model.py}, lines 44--131}]
def round_payload_mb() -> tuple[float, float]:
    #(rx_mb, tx_mb) di un round: modello intero in giu', update stessa taglia in su
    return MODEL_SIZE_MB, MODEL_SIZE_MB


def _payload_energy_j(rx_mb: float, tx_mb: float) -> float:
    return rx_mb * COMM_ENERGY_DL_J_PER_MB + tx_mb * COMM_ENERGY_UL_J_PER_MB


def _payload_time_s(rx_mb: float, tx_mb: float) -> float:
    #RX e TX sequenziali (non sovrapposti), ciascuno con la sua banda
    return rx_mb / COMM_BANDWIDTH_DL_MB_PER_S + tx_mb / COMM_BANDWIDTH_UL_MB_PER_S


def session_energy_j() -> float:
    """Energia fissa di una sessione radio: promozione + coda."""
    return (SESSION_PROMO_POWER_W * SESSION_PROMO_TIME_S
            + SESSION_TAIL_POWER_W * SESSION_TAIL_TIME_S)


def session_time_s() -> float:
    return SESSION_PROMO_TIME_S + SESSION_TAIL_TIME_S


# --------------------------------------------------- scambio del modello
def comm_time_s(profile: DeviceProfile | None = None) -> float:
    #trasferimento giu' e su, piu' le due sessioni radio. `profile` accettato
    #e ignorato: rete unica.
    rx_mb, tx_mb = round_payload_mb()
    return _payload_time_s(rx_mb, tx_mb) + 2 * session_time_s()


def comm_energy_j(profile: DeviceProfile | None = None) -> float:
    rx_mb, tx_mb = round_payload_mb()
    return _payload_energy_j(rx_mb, tx_mb) + 2 * session_energy_j()


def comm_power_w(profile: DeviceProfile | None = None) -> float:
    #potenza media equivalente: serve al modello di batteria, che ragiona in
    #potenza x tempo, cosi' la non-idealita' si applica anche alla comunicazione
    t = comm_time_s()
    return comm_energy_j() / t if t > 0 else 0.0


# ------------------------------------------------ solo download del modello
# [B] 04/11: chi riceve il modello e poi si ritira (opt-out di SAGE-smart)
# scarica il modello ma non carica l'update: una sola sessione radio.
def download_time_s() -> float:
    rx_mb, _ = round_payload_mb()
    return _payload_time_s(rx_mb, 0.0) + session_time_s()


def download_energy_j() -> float:
    rx_mb, _ = round_payload_mb()
    return _payload_energy_j(rx_mb, 0.0) + session_energy_j()


def download_power_w() -> float:
    t = download_time_s()
    return download_energy_j() / t if t > 0 else 0.0


# ------------------------------------------------- messaggi di controllo
# Richiesta di stato o proposta di lavoro, con la risposta del client.
# merged=True: lo scambio cade nella stessa sessione radio del download del
# modello che segue subito dopo (client selezionato), quindi non paga una
# sessione propria. merged=False: sessione isolata, costo pieno.
def ctrl_energy_j(merged: bool = False) -> float:
    half = CTRL_PAYLOAD_MB / 2.0
    return _payload_energy_j(half, half) + (0.0 if merged else session_energy_j())


def ctrl_time_s(merged: bool = False) -> float:
    """Durata della scarica: trasmissione piu', se isolata, la sessione."""
    half = CTRL_PAYLOAD_MB / 2.0
    return _payload_time_s(half, half) + (0.0 if merged else session_time_s())


def ctrl_power_w(merged: bool = False) -> float:
    t = ctrl_time_s(merged)
    return ctrl_energy_j(merged) / t if t > 0 else 0.0


def ctrl_latency_s() -> float:
    """Quanto uno scambio di controllo ritarda il round: promozione + RTT +
    trasmissione. La coda non conta, la radio si spegne in background."""
    half = CTRL_PAYLOAD_MB / 2.0
    return SESSION_PROMO_TIME_S + COMM_RTT_S + _payload_time_s(half, half)
\end{lstlisting}

\subsection{Battery}

The battery of the world (\cf{lin}, \cf{peuk} or \cf{nm}) is a static property
of the phone, known to the server from the enrollment and to the client. A
power $P$ for $\Delta t$ seconds at SoC $s$ changes the SoC by:
\begin{align}
  \text{Ideal:} \quad & \Delta s = -\frac{(P/\eta)\,\Delta t / 3600}{V_0 \, Q_i} , \\
  \text{Peukert:} \quad & I = \frac{P/\eta}{V_0}, \quad
      Q'_i = Q_i \Bigl(\frac{Q_i}{5\,\text{h}}\Bigr)^{n-1}, \quad
      \Delta s = -\frac{\Delta t / 3600}{Q'_i / I^{n}} , \quad n = 1.05 , \\
  \text{Datasheet:} \quad & \frac{P}{\eta} = V\!\Bigl(s, \frac{I}{Q_i}\Bigr) \, I
      \;\text{(solved for } I \text{ by bisection)}, \quad
      \Delta s = -\frac{I \, \Delta t / 3600}{Q_i \, q(I/Q_i)} .
\end{align}
In the Datasheet model $V(s, \rho)$ is interpolated linearly in $s$ on the three
discharge curves of the datasheet (0.2C, 0.5C, 1C) and linearly in the C-rate
$\rho$; below 0.2C it is the 0.2C curve, above 1C it is extrapolated (at least
0.5 V). The usable capacity is $q(\rho) = 1.000, 0.962, 0.919$ of the nominal
one at 0.2C, 0.5C and 1C (constant below 0.2C, extrapolated above, at least
0.3). The bisection runs 60 iterations between $10^{-9}$ A and $50 Q_i$. The
world clamps the SoC to $[0, 1]$; a phone that reaches 0 is dead for the rest
of the run. A phone in charge is connected to the grid: its battery does not
discharge, and the energy is counted separately as grid energy.

\begin{lstlisting}[style=py,firstnumber=71,caption={\texttt{tesifl/device/world\_state.py}, lines 71--87}]
def battery_delta_soc(battery: str, power_w: float, dt_s: float,
                      profile: DeviceProfile, soc: float) -> float:
    """Variazione di SoC (<= 0) della batteria `battery` per power_w x dt_s.

    lin   SoC -= P/eta * dt / (V_nom * C_nom)                (energy_wh_linear)
    peuk  SoC -= dt / t_pieno,  t_pieno = C' / I^n           (battery_model_peukert)
    nm    SoC -= I * dt / Q(I), P = V(SoC, I) * I            (battery_model)
    Solo nm dipende dal SoC corrente: la tensione cala al calare della carica.
    """
    if battery == "nm":
        return delta_soc(power_w, dt_s, profile, soc)
    if battery == "peuk":
        return peukert_delta_soc(power_w, dt_s, profile)
    if battery == "lin":
        cap_wh = profile.battery_capacity_mah / 1000.0 * V_NOMINAL
        return -energy_wh_linear(power_w, dt_s) / cap_wh
    raise ValueError(f"batteria '{battery}': una fra {BATTERIES}")
\end{lstlisting}

\begin{lstlisting}[style=py,firstnumber=67,caption={\texttt{tesifl/device/battery\_model.py}, lines 67--90}]
_SOC = np.array([0.00, 0.01, 0.02, 0.03, 0.04, 0.05, 0.075, 0.10, 0.15,
                 0.20, 0.25, 0.30, 0.35, 0.40, 0.45, 0.50, 0.55, 0.60,
                 0.65, 0.70, 0.75, 0.80, 0.85, 0.90, 0.95, 0.98, 1.00])

# V a 0.2C -- la curva di gran lunga piu' usata: training ~0.2C, idle ~0.01C
_V_02C = np.array([2.50, 3.00, 3.15, 3.26, 3.33, 3.38, 3.44, 3.46, 3.51,
                   3.54, 3.57, 3.59, 3.61, 3.64, 3.66, 3.69, 3.71, 3.75,
                   3.80, 3.85, 3.90, 3.95, 4.03, 4.09, 4.17, 4.22, 4.24])
# V a 0.5C
_V_05C = np.array([2.50, 2.80, 2.94, 3.04, 3.11, 3.16, 3.24, 3.29, 3.36,
                   3.41, 3.44, 3.47, 3.49, 3.51, 3.53, 3.55, 3.59, 3.62,
                   3.66, 3.71, 3.76, 3.81, 3.89, 3.96, 4.03, 4.08, 4.12])
# V a 1C
_V_1C = np.array([2.50, 2.68, 2.78, 2.86, 2.91, 2.96, 3.04, 3.10, 3.18,
                  3.23, 3.27, 3.30, 3.33, 3.35, 3.38, 3.40, 3.43, 3.47,
                  3.50, 3.55, 3.59, 3.65, 3.72, 3.78, 3.84, 3.89, 3.95])

_CRATES = np.array([0.2, 0.5, 1.0])
_VTAB = np.vstack([_V_02C, _V_05C, _V_1C])          # (3, 27)

# capacita' utilizzabile come FRAZIONE della nominale, dagli endpoint
_QFRAC = np.array([2.60, 2.50, 2.39]) / 2.60         # 1.000, 0.962, 0.919

_BISECT_ITERS = 60
\end{lstlisting}

\begin{lstlisting}[style=py,firstnumber=93,caption={\texttt{tesifl/device/battery\_model.py}, lines 93--159}]
def _voltage(soc: float, crate: float) -> float:
    """V(SoC, C-rate), bilineare sulle tre curve.

    [B] fuori dal range 0.2-1C si estrapola LINEARMENTE in C-rate invece di
    saturare: a correnti molto basse (idle, ~0.001C) la tensione satura da
    sola sulla curva 0.2C, che e' fisicamente corretto, mentre a correnti alte
    saturare sottostimerebbe la caduta e falserebbe proprio il caso peggiore.
    """
    soc = float(np.clip(soc, 0.0, 1.0))
    v_at = np.array([np.interp(soc, _SOC, _VTAB[i]) for i in range(3)])
    if crate <= _CRATES[0]:
        # sotto 0.2C la caduta ohmica e' gia' trascurabile: non estrapolare
        # verso l'alto, la tensione a vuoto e' il limite fisico
        return float(v_at[0])
    if crate >= _CRATES[-1]:
        slope = (v_at[2] - v_at[1]) / (_CRATES[2] - _CRATES[1])
        return float(max(0.5, v_at[2] + slope * (crate - _CRATES[-1])))
    return float(np.interp(crate, _CRATES, v_at))


def _q_fraction(crate: float) -> float:
    """Capacita' utilizzabile / nominale, in funzione del C-rate."""
    if crate <= _CRATES[0]:
        return float(_QFRAC[0])
    if crate >= _CRATES[-1]:
        slope = (_QFRAC[2] - _QFRAC[1]) / (_CRATES[2] - _CRATES[1])
        return float(max(0.3, _QFRAC[2] + slope * (crate - _CRATES[-1])))
    return float(np.interp(crate, _CRATES, _QFRAC))


def solve_current_a(power_w: float, soc: float, profile: DeviceProfile) -> float:
    """Risolve P = V(SoC, I) * I per bisezione. Ritorna I in ampere.

    [B] serve un metodo iterativo perche' V dipende da I: piu' corrente, meno
    tensione, quindi per la stessa potenza serve ancora piu' corrente. Il
    prodotto V*I e' monotono crescente in I finche' la caduta resta moderata,
    il che rende la bisezione sicura e sufficiente -- Newton converge prima ma
    puo' uscire dal dominio tabulato.
    """
    if power_w <= 0.0:
        return 0.0
    q_nom_ah = profile.battery_capacity_mah / 1000.0
    p_batt = power_w / DCDC_EFFICIENCY        # la conversione avviene a monte
    lo, hi = 1e-9, 50.0 * q_nom_ah            # fino a 50C: limite di sicurezza
    for _ in range(_BISECT_ITERS):
        mid = 0.5 * (lo + hi)
        if _voltage(soc, mid / q_nom_ah) * mid < p_batt:
            lo = mid
        else:
            hi = mid
    return 0.5 * (lo + hi)


def delta_soc(power_w: float, dt_s: float, profile: DeviceProfile,
              soc: float = 1.0) -> float:
    """Variazione di SoC (<= 0) consumando power_w per dt_s secondi.

    [B] `soc` ha default 1.0 SOLO per compatibilita' con le chiamate che non
    lo passano ancora. Passarlo cambia il risultato -- e' tutto il punto del
    modello -- quindi ogni chiamante dovrebbe fornire il SoC corrente.
    """
    if power_w <= 0.0 or dt_s <= 0.0:
        return 0.0
    q_nom_ah = profile.battery_capacity_mah / 1000.0
    i_a = solve_current_a(power_w, soc, profile)
    q_eff_ah = q_nom_ah * _q_fraction(i_a / q_nom_ah)
    return -(i_a * (dt_s / 3600.0)) / q_eff_ah
\end{lstlisting}

\begin{lstlisting}[style=py,firstnumber=173,caption={\texttt{tesifl/device/battery\_model.py}, lines 173--189}]
def energy_wh_linear(power_w: float, dt_s: float) -> float:
    """Energia al carico, P/eta * dt, senza alcun modello di batteria.

    E' anche il consumo della batteria LINEARE (mondo "lin"): li' il SoC
    scende di P/eta * dt / (V_nom * C_nom). Negli altri mondi resta la
    colonna energy_lin_wh del CSV: l'energia che il device ha usato, a
    prescindere da quanta carica la batteria abbia pagato per erogarla.
    """
    if power_w <= 0.0 or dt_s <= 0.0:
        return 0.0
    return (power_w / DCDC_EFFICIENCY) * dt_s / 3600.0


def energy_wh_from_delta_soc(delta: float, profile: DeviceProfile) -> float:
    """Energia estratta DALLA BATTERIA per una data variazione di SoC."""
    q_nom_ah = profile.battery_capacity_mah / 1000.0
    return V_NOMINAL * (-delta * q_nom_ah)
\end{lstlisting}

\begin{lstlisting}[style=py,firstnumber=29,caption={\texttt{tesifl/device/battery\_model\_peukert.py}, lines 29--49}]
def load_power_to_battery_current(power_w: float) -> float:
    """P_batt = P_load / eta_dcdc ; I_batt = P_batt / V_nom."""
    if power_w <= 0:
        return 0.0
    return (power_w / DCDC_EFFICIENCY) / V_NOMINAL


def peukert_capacity_ah(profile: DeviceProfile) -> float:
    """C' = C * (C/R)^(n-1), la capacita' effettiva alla corrente nominale."""
    c_rated_ah = profile.battery_capacity_mah / 1000.0
    i_rated_a = c_rated_ah / RATED_R_HOURS
    return c_rated_ah * (i_rated_a ** (profile.peukert_n - 1.0))


def delta_soc(power_w: float, dt_s: float, profile: DeviceProfile) -> float:
    """Variazione di SoC (<= 0) consumando power_w per dt_s secondi."""
    if power_w <= 0 or dt_s <= 0:
        return 0.0
    current_a = load_power_to_battery_current(power_w)
    t_full_h = peukert_capacity_ah(profile) / (current_a ** profile.peukert_n)
    return -((dt_s / 3600.0) / t_full_h)
\end{lstlisting}

\subsection{Workload of the other apps}

Each phone follows one of three Markov chains over the states idle, light and
heavy ($w = 0$, 0.10, 0.30). The chains are assigned in equal thirds with a
permutation independent of the tier (seed of the run + 7919) and make one
transition per round. Their stationary means are 0.0297, 0.1191 and 0.2055;
their mean, $\bar{w} = 0.1181$, is the workload that the server assumes for a
phone that has never reported one. The heaviest state, $w_{\max} = 0.30$, is
only a fallback value for a reply without workload.

\begin{lstlisting}[style=py,firstnumber=41,caption={\texttt{tesifl/device/workload.py}, lines 41--41}]
W_MAX = 0.95
\end{lstlisting}

\begin{lstlisting}[style=py,firstnumber=44,caption={\texttt{tesifl/device/workload.py}, lines 44--51}]
def scaled_util(scale: float) -> np.ndarray:
    """Utilizzazione dei tre stati moltiplicata per `scale`, con tetto W_MAX.

    A scale = 1 i valori sono quelli di experiment.toml, identici bit per bit.
    """
    if scale <= 0.0:
        raise ValueError(f"workload scale {scale}: deve essere > 0")
    return np.minimum(W_MAX, STATE_UTIL * float(scale))
\end{lstlisting}

\begin{lstlisting}[style=py,firstnumber=67,caption={\texttt{tesifl/device/workload.py}, lines 67--112}]
    def __init__(self, n_clients: int, seed: int, enabled: bool = True,
                 scale: float = 1.0) -> None:
        self.enabled = bool(enabled)
        self.scale = float(scale)
        self._util = scaled_util(self.scale)
        self.n_clients = int(n_clients)
        self._rng = np.random.default_rng(seed)
        self._chains = {p: np.asarray(WORKLOAD_CHAINS[p], dtype=float)
                        for p in PROFILE_NAMES}
        for p, m in self._chains.items():
            if not np.allclose(m.sum(axis=1), 1.0):
                raise ValueError(f"la catena '{p}' ha righe che non sommano a 1")

        # profili in parti uguali, assegnati con una permutazione: senza,
        # il profilo sarebbe una funzione del cid e quindi correlato alla
        # partizione dei dati, che pure dipende dal cid.
        self._profile = self._assign_equal(PROFILE_NAMES)
        # tutti partono da idle: lo stato iniziale conta poco, le catene
        # convergono alla stazionaria in poche decine di round.
        self._state = {cid: 0 for cid in range(self.n_clients)}

    def _assign_equal(self, names: list[str]) -> dict[int, str]:
        base = [names[i % len(names)] for i in range(self.n_clients)]
        order = self._rng.permutation(self.n_clients)
        return {int(cid): base[i] for i, cid in enumerate(order)}

    # ------------------------------------------------------------ evoluzione
    def step(self) -> None:
        """Una transizione per ogni client. Da chiamare a inizio round."""
        if not self.enabled:
            return
        for cid in range(self.n_clients):
            row = self._chains[self._profile[cid]][self._state[cid]]
            self._state[cid] = int(self._rng.choice(len(row), p=row))

    # -------------------------------------------------------------- lettura
    def utilization(self, cid: int) -> float:
        """w in [0, 1): frazione di calcolo occupata da altre applicazioni."""
        if not self.enabled:
            return 0.0
        return float(self._util[self._state[cid]])

    def max_utilization(self) -> float:
        """w dello stato piu' pesante a questo scale: il caso peggiore che
        SAGE-smart assume finche' non conosce il workload di un client."""
        return float(self._util.max())
\end{lstlisting}

\begin{lstlisting}[style=py,firstnumber=132,caption={\texttt{tesifl/device/workload.py}, lines 132--147}]
def stationary_utilization(scale: float = 1.0) -> dict[str, float]:
    """Utilizzazione media attesa di ogni profilo, dall'autovettore della catena.

    Serve a verificare che le matrici diano le medie che ci si aspetta, senza
    doverle simulare: e' il primo dei controlli di validazione.
    """
    util = scaled_util(scale)
    out = {}
    for name, m in WORKLOAD_CHAINS.items():
        m = np.asarray(m, dtype=float)
        vals, vecs = np.linalg.eig(m.T)
        i = int(np.argmin(np.abs(vals - 1.0)))
        pi = np.real(vecs[:, i])
        pi = pi / pi.sum()
        out[name] = float(pi @ util)
    return out
\end{lstlisting}

\subsection{Initial SoC, recharge and idle of the world}

The initial SoC is drawn from a normal distribution with mean 0.75 and standard
deviation 0.15, rejecting the samples outside $[0.30, 1.00]$. At the start of
every round a phone that is alive and not charging is plugged in with
probability
\begin{equation}
  p^{\mathrm{plug}}_i = 1 - \exp\!\bigl(-\lambda \, (1 - s_i) \, \Delta t_{r-1} / 3600\bigr) ,
  \label{eq:plug}
\end{equation}
with $\Delta t_{r-1}$ the duration of the previous round in seconds (0 at round
1). A charging phone gains $\gamma \, t_r / 3600$ of SoC per round of duration
$t_r$ and leaves the charge at 1.0. The round lasts as long as the slowest
selected phone plus the control overhead; every alive phone is idle, at the
power of Equation~\eqref{eq:pidle} with its current workload, for the part of
the round in which it is not training.

\begin{lstlisting}[style=py,firstnumber=191,caption={\texttt{tesifl/device/world\_state.py}, lines 191--219}]
    def _sample_initial_soc(self) -> float:
        """SoC iniziale: normale troncata, spostata verso la carica piena.

        [B] l'uniforme su [0.3, 1.0] metteva il 29% dei client sotto il 50%
        di carica gia' al round 1, cioe' una popolazione mediamente scarica
        prima ancora di iniziare. Con mu=0.75 quella quota scende al 5%: resta
        una coda bassa -- che serve, perche' e' cio' che gli algoritmi
        battery-aware devono saper distinguere -- ma il caso tipico diventa un
        device carico, che e' la situazione realistica.

        [B] RIFIUTO IN CICLO, NON CLIP. Un clip agli estremi accumulerebbe
        massa su 0.30 e 1.00, creando due picchi artificiali proprio dove la
        soglia di riserva e la carica piena contano di piu'.

        [B] usa self._rng, quindi la popolazione iniziale dipende dal seed
        come prima: seed diversi danno popolazioni diverse, lo stesso seed
        la stessa popolazione.
        """
        lo, hi = INITIAL_SOC_RANGE
        for _ in range(10000):
            x = float(self._rng.normal(INITIAL_SOC_MEAN, INITIAL_SOC_STD))
            if lo <= x <= hi:
                return x
        # irraggiungibile con parametri sensati; se mu cade fuori dal
        # troncamento e' meglio accorgersene che restare in ciclo
        raise RuntimeError(
            f"SoC iniziale: nessun campione in [{lo}, {hi}] con "
            f"mu={INITIAL_SOC_MEAN}, sigma={INITIAL_SOC_STD}"
        )
\end{lstlisting}

\begin{lstlisting}[style=py,firstnumber=298,caption={\texttt{tesifl/device/world\_state.py}, lines 298--320}]
    def begin_round(self) -> None:
        """Inizio round: avanza il workload e decide chi va in carica.

        [B] la probabilita' cresce al calare del SoC: e' il comportamento
        dell'utente, che mette in carica quando vede la batteria bassa. E' un
        tasso ORARIO applicato alla durata dt del round appena finito,
            p = 1 - exp(-rate_per_hour * (1 - SoC) * dt_ore),
        cosi' algoritmi con round di durata diversa vedono la stessa ricarica
        per ora di tempo reale (prima p era fissa per round). Al primo round
        dt = 0: nessuno parte in carica. Un device gia' in carica non viene
        riestratto: esce dalla carica solo al 100%.
        """
        self.workload.step()
        if not self._recharge_enabled:
            return
        dt_h = self._last_round_s / 3600.0
        for cid in self._profiles:
            if self._failed[cid] or self._charging[cid]:
                continue
            p = 1.0 - np.exp(-RECHARGE_RATE_PER_H * (1.0 - self._soc[cid]) * dt_h)
            if self._rng.random() < p:
                self._charging[cid] = True
                self._n_recharges[cid] += 1
\end{lstlisting}

\begin{lstlisting}[style=py,firstnumber=332,caption={\texttt{tesifl/device/world\_state.py}, lines 332--352}]
    def round_duration_s(
        self, cid: int, epochs: int, batch_size: int, dataset_size_local: int,
        include_workload: bool = True,
    ) -> float:
        """Durata del training locale.

        include_workload=False da' la durata NOMINALE, che dipende dal solo
        tier. Serve a ESCS, che nel paper usa la latenza di profilo come dato
        statico noto al server: il workload e' proprio cio' che il server non
        puo' conoscere in anticipo.
        """
        p = self.effective_profile(cid) if include_workload else self._profiles[cid]
        return training_time_s(p, epochs, dataset_size_local, batch_size)

    def probe_duration_s(self, cid: int, n_samples: int) -> float:
        """Forward di sola valutazione su n_samples, col workload corrente.

        E' la stima della loss locale che il client di sage_smart fa prima
        di decidere (E, B): costa poco, ma non e' gratis e va addebitata.
        """
        return probe_time_s(self.effective_profile(cid), n_samples)
\end{lstlisting}

\begin{lstlisting}[style=py,firstnumber=362,caption={\texttt{tesifl/device/world\_state.py}, lines 362--423}]
    def _drain(self, cid: int, power_w: float, dt_s: float) -> tuple[float, float]:
        """Applica power_w per dt_s con la batteria del mondo.

        Ritorna (Wh di carica pagati dalla batteria, Wh al carico P*dt).
        """
        if self._failed[cid]:
            return 0.0, 0.0
        if self._charging[cid]:
            # collegato alla rete: la batteria non si scarica. L'energia c'e'
            # lo stesso e la contabilizziamo a parte.
            self._energy_grid_wh[cid] += (power_w / DCDC_EFFICIENCY) * dt_s / 3600.0
            return 0.0, 0.0
        profile = self._profiles[cid]
        soc = self._soc[cid]
        d_req = battery_delta_soc(self.battery, power_w, dt_s, profile, soc)
        new_soc = max(0.0, min(1.0, soc + d_req))
        # [B] energia dal delta EFFETTIVO, non da quello richiesto: il SoC e'
        # clampato a 0 ma il delta no, quindi un client che muore nel round
        # verrebbe addebitato per carica che non aveva.
        d = new_soc - soc
        self._soc[cid] = new_soc
        if new_soc <= 0.0:
            self._failed[cid] = True
        # l'energia al carico va scalata con la stessa frazione effettiva: se
        # il client muore a meta' round, ha lavorato solo per meta' round
        frac = d / d_req if d_req < 0.0 else 0.0
        e_lin = energy_wh_linear(power_w, dt_s) * frac
        self._energy_lin_wh[cid] += e_lin
        return energy_wh_from_delta_soc(d, profile), e_lin

    def apply_round(self, cid: int, dt_s: float, batch_size: int) -> tuple[float, float]:
        e, e_lin = self._drain(cid, self.training_power_w(cid, batch_size), dt_s)
        self._energy_wh[cid] += e
        return e, e_lin

    def apply_control(self, cid: int, merged: bool = False) -> tuple[float, float]:
        """Uno scambio di controllo (richiesta + risposta) con il server.

        merged=True se lo scambio cade nella sessione radio del download del
        modello che segue (client selezionato): paga il solo payload. Vedi
        communication_model.
        """
        e, e_lin = self._drain(cid, ctrl_power_w(merged), ctrl_time_s(merged))
        self._energy_ctrl_wh[cid] += e
        return e, e_lin

    def apply_communication(self, cid: int) -> tuple[float, float]:
        e, e_lin = self._drain(cid, comm_power_w(), comm_time_s())
        self._energy_comm_wh[cid] += e
        return e, e_lin

    def apply_download(self, cid: int) -> tuple[float, float]:
        """Solo il download del modello, con la sua sessione radio.

        [B] 04/11: e' cio' che paga un client di SAGE-smart che riceve il
        modello e poi si ritira (opt-out): prima pagava lo scambio intero
        (~181 J con ResNet-18), anche l'upload che non fa (~46 J il solo
        download). La risposta di ritiro e' un messaggio di controllo a parte.
        """
        e, e_lin = self._drain(cid, download_power_w(), download_time_s())
        self._energy_comm_wh[cid] += e
        return e, e_lin
\end{lstlisting}

\begin{lstlisting}[style=py,firstnumber=425,caption={\texttt{tesifl/device/world\_state.py}, lines 425--475}]
    def idle_power_w(self, cid: int) -> float:
        """P_idle + w * (P_train - P_idle) del tier, vedi idle_power_w()."""
        return idle_power_w(self._profiles[cid], self.workload.utilization(cid))

    def apply_idle(self, durations: dict[int, float]) -> tuple[float, float]:
        """Idle di fine round, modello ESCS Sez. 4.2.

        durations: cid -> durata del training, per i soli SELEZIONATI.
        Il round dura quanto il piu' lento fra loro; ogni client sta in idle
        per la differenza fra quella durata e il proprio training (zero per il
        piu' lento, tutta la durata per chi non e' stato selezionato).

        [B] va chiamata DOPO apply_round, altrimenti un client gia' a zero
        pagherebbe l'idle prima del training e l'ordine cambierebbe chi muore.
        """
        if not self._idle_enabled:
            return 0.0, 0.0
        t_round = max(durations.values()) if durations else 0.0
        if t_round <= 0.0:
            return 0.0, 0.0
        self._last_round_s = t_round
        self._charge_step(t_round)
        tot_e = tot_lin = 0.0
        for cid in self._profiles:
            if self._failed[cid]:
                continue
            idle_s = t_round - float(durations.get(cid, 0.0))
            if idle_s <= 0.0:
                continue
            e, lin = self._drain(cid, self.idle_power_w(cid), idle_s)
            self._energy_idle_wh[cid] += e
            tot_e += e
            tot_lin += lin
        return tot_e, tot_lin

    def _charge_step(self, t_round: float) -> None:
        """Carica lineare a C-rate costante per la durata del round.

        [B] lineare per scelta: la curva reale e' CC-CV e rallenta sopra l'80%,
        ma modellarla aggiungerebbe due parametri per un effetto che qui non
        cambia nulla. La carica finisce al 100% e il device torna normale.
        """
        if not self._recharge_enabled:
            return
        d_soc = CHARGE_C_RATE * t_round / 3600.0
        for cid in self._profiles:
            if not self._charging[cid]:
                continue
            self._soc[cid] = min(1.0, self._soc[cid] + d_soc)
            if self._soc[cid] >= 1.0:
                self._charging[cid] = False
\end{lstlisting}

\newpage
%======================================================================
\section{Structure of a round}
\label{sec:round}

\cf{PhysicalFedAvg.configure\_train} and \cf{aggregate\_train} define the order
of the operations. For SAGE-smart one round $r$ is:
\begin{enumerate}
  \item $\mathcal{A}_r$ = phones not depleted. Phones in charge stay available
        (\cf{recharge\_available = true}). No broadcast.
  \item \cf{\_select\_clients} (Sections~\ref{sec:enroll}--\ref{sec:fresh}):
        enrollment at round 1, forecast, ranking, number of slots, proposals,
        fresh choice, energy of the control messages, envelopes. The phones see
        the proposals with their state of the previous round.
  \item \cf{world.begin\_round()}: one workload transition per phone and the
        recharge draws of Equation~\eqref{eq:plug}.
  \item The model, the envelope and the sensor readings of this round (after
        step 3) go to the phones of $S_r$ (late binding).
  \item Each phone decides $(E, B)$, trains and replies
        (Section~\ref{sec:client}).
  \item \cf{aggregate\_train} (Section~\ref{sec:aggregate}): the replies are
        sorted by phone index; for each one the world charges training and
        communication and the strategy updates its state. Then the round
        duration $t_r = \max_{i \in S_r} t_i + t^{\mathrm{ovh}}_r$ is fixed,
        the charging phones recharge, every alive phone pays the idle, and the
        models are averaged.
  \item The global model is evaluated on the test set.
\end{enumerate}

\begin{lstlisting}[style=py,firstnumber=110,caption={\texttt{tesifl/strategies/fedavg.py}, lines 110--168}]
    def configure_train(
        self, server_round: int, arrays: ArrayRecord, config: ConfigRecord, grid: Grid
    ) -> list[Message]:
        if self._exhausted:
            return []
        self._ensure_mapping(grid)
        self.world.tick_all()
        available = [
            nid for nid, cid in self._node_to_cid.items()
            if self._is_available(cid)
        ]
        self._n_alive = len(available)
        if self.broadcast:
            available = self._broadcast(available, server_round)
        k = min(self.k_per_round, len(available))
        selected = self._select_clients(available, k, server_round)
        if not selected:
            self._exhausted = True
            log(INFO, "round %s: selezione vuota su %s disponibili, run terminata",
                server_round, len(available))
            return []
        if self.broadcast:
            self._charge_broadcast(selected)

        # [B] il workload avanza qui, prima che il server calcoli le durate e
        # selezioni: cosi' la selezione di questo round vede lo stato corrente,
        # non quello del round precedente.
        self.world.begin_round()

        self._global_arrays = arrays

        base = dict(config)
        base["server-round"] = server_round
        base["epochs"] = self.epochs
        base["batch-size"] = self.batch_size

        self._epochs_sent.clear()
        self._batch_sent.clear()
        messages: list[Message] = []
        for nid in selected:
            cid = int(self._node_to_cid[int(nid)])
            cfg = ConfigRecord(dict(base))
            cfg["cid"] = cid
            self._per_client_config(cid, cfg, server_round)
            self._epochs_sent[cid] = int(cfg["epochs"])
            self._batch_sent[cid] = int(cfg["batch-size"])
            rec = RecordDict(
                {self.arrayrecord_key: arrays, self.configrecord_key: cfg}
            )
            dev = self._device_record(cid)
            if dev is not None:
                rec["device"] = dev
            messages.extend(
                self._construct_messages(rec, [int(nid)], MessageType.TRAIN)
            )

        log(INFO, "round %s: %s/%s disponibili, %s selezionati",
            server_round, len(available), len(self._node_to_cid), len(selected))
        return messages
\end{lstlisting}

\begin{lstlisting}[style=py,firstnumber=297,caption={\texttt{tesifl/strategies/fedavg.py}, lines 297--361}]
    def aggregate_train(
        self, server_round: int, replies: list[Message]
    ) -> tuple[ArrayRecord | None, MetricRecord | None]:
        # nessuna reply: niente da aggregare, e super() dividerebbe per zero
        self._last_round_s = float(self._round_overhead_s())
        if not replies:
            return None, None

        replies = sorted(
            replies,
            key=lambda r: self._node_to_cid.get(int(r.metadata.src_node_id), 1 << 62),
        )

        durations: dict[int, float] = {}
        for reply in replies:
            if not reply.has_content():
                continue
            cid = self._node_to_cid[reply.metadata.src_node_id]
            n_ex = int(self._metric(reply, "num-examples", "num_examples"))
            dt, b = self._work_done(cid, n_ex, reply)
            durations[cid] = dt
            e_tr, l_tr = self.world.apply_round(cid, dt_s=dt, batch_size=b)
            e_cm, l_cm = self._charge_communication(cid, reply)
            if n_ex > 0:
                self._participations[cid] = self._participations.get(cid, 0) + 1
            self._on_round_result(cid, n_ex, dt, float(e_tr + e_cm),
                                  float(l_tr + l_cm), reply, server_round)

        # [B] idle DOPO il training: il round dura quanto il selezionato piu'
        # lento, e ogni client paga la differenza (modello ESCS Sez. 4.2). Chi
        # non e' stato selezionato paga l'intera durata. L'ordine conta: se
        # l'idle venisse prima, un client al limite morirebbe di idle invece
        # che di training e il conteggio cambierebbe.
        # [B] il ritardo prima del training (messaggi di controllo di
        # sage_smart) allunga il round: in idle lo pagano tutti tranne i
        # selezionati, che in quel tempo stanno parlando col server.
        overhead = float(self._round_overhead_s())
        self._last_round_s = max(durations.values(), default=0.0) + overhead
        self._sim_time_s += self._last_round_s
        self.world.apply_idle({c: d + overhead for c, d in durations.items()})
        self._dev_soc[int(server_round)] = self.world.soc_vector()

        s = self.world.stats()
        self._round_stats[int(server_round)] = dict(s)
        self._round_stats[int(server_round)].update({
            "round_s": float(self._last_round_s),
            "sim_time_s": float(self._sim_time_s),
        })
        if self.broadcast:
            # [B] l'energia del broadcast e' gia' dentro energy_ctrl_wh (e
            # quindi nel totale): qui la si riporta a parte per leggerla
            self._round_stats[int(server_round)].update({
                **self._bcast_round,
                "bcast_overhead_s": float(self._bcast_overhead_s),
                "energy_bcast_wh": float(self._bcast_energy_wh),
            })
        log(INFO, "round %s: %.3f Wh (train %.3f + comm %.3f) | SoC medio %.2f "
            "| in ricarica %s | falliti %s",
            server_round, s["total_energy_wh"], s["energy_train_wh"],
            s["energy_comm_wh"], s["mean_soc"], s["n_recharging"], s["n_failed"])

        replies = self._aggregatable(replies)
        if not replies:
            return None, None
        return super().aggregate_train(server_round, replies)
\end{lstlisting}

\begin{lstlisting}[style=py,firstnumber=333,caption={\texttt{tesifl/strategies/sage\_smart.py}, lines 333--341}]
    def configure_train(self, server_round, arrays, config, grid):
        # la griglia serve alla negoziazione per i messaggi di controllo, che
        # partono dentro _select_clients prima dei messaggi di training
        self._grid = grid
        return super().configure_train(server_round, arrays, config, grid)

    def _round_overhead_s(self) -> float:
        # negoziazione + broadcast dello stato (solo con noneg, vedi PhysicalFedAvg)
        return self._overhead_s + super()._round_overhead_s()
\end{lstlisting}

\newpage
%======================================================================
\section{The strategy object}
\label{sec:init}

The constructor fixes the components that are on (all of them, unless an
ablation switches them off), the slots, the batch sizes of the envelope, the
fallback workload $w_{\max}$ and the enrollment workload $\bar{w}$, and creates
the state that the server keeps for every phone: the last report (SoC,
workload, charging, time), the known state, $k_i$, the class histograms and
the last round of contact. Without broadcast (\cf{self.broadcast} false, as set
by \cf{server\_app} whenever the negotiation is on) the fraction $p_{\mathrm{nr}}$
of the contacted phones does not answer (\cf{proposal\_nonresp}).

\begin{lstlisting}[style=py,firstnumber=212,caption={\texttt{tesifl/strategies/sage\_smart.py}, lines 212--312}]
    def __init__(
        self,
        *args,
        stale_weight: float = 0.3,
        stale_max: int = 25,
        soc_min: float = 0.20,
        epochs_cap_mult: int = 2,
        epochs_min: int = 2,
        partition_sizes: dict[int, int] | None = None,
        deadline_mult: float = 1.0,
        probe_samples: int = 256,
        max_waves: int = 2,
        query_timeout_s: float = 600.0,
        calib_alpha: float = 0.3,
        batch_min: int = 16,
        batch_max: int = 64,
        forecast: bool = True,
        negotiation: bool = True,
        partial_work: bool = True,
        shortlist: float = 1.5,
        load_tau_s: float = 450.0,
        **kwargs,
    ) -> None:
        super().__init__(*args, **kwargs)
        if shortlist < 1.0:
            raise ValueError(f"shortlist >= 1, non {shortlist}")
        self.forecast = bool(forecast)
        self.negotiation = bool(negotiation)
        self.partial_work = bool(partial_work)
        self.shortlist = float(shortlist)
        self.load_tau_s = float(load_tau_s)
        # [B] senza broadcast nessuno ha verificato chi e' raggiungibile: la
        # frazione broadcast_nonresp dei contattati non risponde (proposta o
        # modello). Il broadcast lo spegne server_app quando c'e' la
        # negoziazione; resta solo nell'ablazione senza negoziazione (noneg).
        self.proposal_nonresp = not self.broadcast
        self.stale_w = float(stale_weight)
        self.stale_max = int(stale_max)
        self.soc_min = float(soc_min)
        if not 1 <= int(batch_min) <= int(batch_max):
            raise ValueError(f"batch: serve 1 <= min <= max, non {batch_min}, {batch_max}")
        choices, b = [], int(batch_min)
        while b <= int(batch_max):
            choices.append(b)
            b *= 2
        self.batch_choices = tuple(choices)
        self.epochs_cap = int(epochs_cap_mult) * int(self.epochs)
        self.epochs_min = int(epochs_min)
        self.deadline_mult = float(deadline_mult)
        self.probe_samples = int(probe_samples)
        self.max_waves = max(1, int(max_waves))
        self.query_timeout_s = float(query_timeout_s)
        self.calib_alpha = float(calib_alpha)
        # Round 0: taglia dei dataset. Se manca, la si impara dalle risposte.
        self._n: dict[int, int] = {int(c): int(n) for c, n in
                                   (partition_sizes or {}).items()}
        # ripiego se una risposta non riportasse il workload: il caso peggiore
        # (segue workload-scale: a x1 e' 0.30)
        self.w_busy = self.world.workload.max_utilization()
        # carico assunto all'iscrizione: la media della POPOLAZIONE (vedi la
        # docstring della classe, round 0)
        wl = self.world.workload
        self.w_enroll = (float(np.mean(list(stationary_utilization(wl.scale).values())))
                         if wl.enabled else 0.0)

        self._last_seen: dict[int, int] = {}       # cid -> ultimo round scelto
        self._plan: dict[int, Envelope] = {}       # cid -> envelope del round

        # stato che il server tiene per ogni client: l'ultimo REPORT (cosa il
        # client ha detto e quando) e lo stato NOTO (previsione o report fresco)
        self._subscribed = False
        self._deadline_s: float = math.inf
        self._clock_s = 0.0                        # tempo simulato, inizio round
        self._soc_rep: dict[int, float] = {}
        self._t_rep: dict[int, float] = {}
        self._w_rep: dict[int, float] = {}
        self._w_sum: dict[int, float] = {}         # workload riportati, per la media
        self._w_cnt: dict[int, int] = {}
        self._chg_rep: dict[int, bool] = {}
        self._soc_known: dict[int, float] = {}
        self._chg_known: dict[int, bool] = {}
        self._w_known: dict[int, float] = {}
        self._k_cal: dict[int, float] = {}
        self._loss_known: dict[int, float] = {}
        # per round: vincolo che ha fissato E_hi, risposte dei client e costi
        # della negoziazione
        self._bound: dict[int, str] = {}
        self._reports: dict[int, dict] = {}
        self._grid = None
        self._queried: set[int] = set()
        self._overhead_s = 0.0
        self._waves = 0
        self._n_queried = 0
        self._n_refused = 0
        self._info_err = float("nan")
        # rinviati, non raggiunti, attesa del timeout per chi non risponde,
        # errore sul workload noto (verifica)
        self._n_deferred = 0
        self._n_noreply = 0
        self._extra_wait_s = 0.0
        self._w_err = float("nan")
\end{lstlisting}

\begin{lstlisting}[style=py,firstnumber=314,caption={\texttt{tesifl/strategies/sage\_smart.py}, lines 314--320}]
    @property
    def _protocol_cols(self) -> bool:
        """Previsione, negoziazione o niente broadcast: colonne in piu' nel CSV
        (rinviati, non raggiunti, errore sul workload). Spente solo nei
        gradini delle cumulate senza stima e senza negoziazione, che cosi'
        hanno le colonne di prima."""
        return self.forecast or self.negotiation or self.proposal_nonresp
\end{lstlisting}

\newpage
%======================================================================
\section{Round 0: enrollment}
\label{sec:enroll}

At the start of round 1, before the first selection, the server registers every
phone with a first report at time 0: the true initial SoC, workload $\bar{w}$,
not charging. This report is marked as not measured, so it does not enter the
mean of the workloads reported by the phone. $k_i = 1$. In the simulator the
enrollment is not a message: it has no energy or time cost. The static data
(tier, throughput, training power, capacity, battery model) are read from the
profile of the phone; $n_i$ comes from \cf{partition\_sizes}, which
reproduces the partition of the clients.

The deadline is the median, over the phones, of the nominal time of the
standard work ($E_0 = 5$, $B_0 = 128$, no workload), by
Equation~\eqref{eq:ttrain}:
\begin{equation}
  T = \texttt{deadline\_mult} \cdot \operatorname{median}_{i \in \mathcal{N}}
      t^{\mathrm{tr}}_i(E_0, n_i, B_0)\big|_{w = 0} .
\end{equation}
It stays fixed for the whole run. With seed 2026, $T = 132.39$ s at $\beta =
0.1$, 130.53 s at $\beta = 0.5$ and 127.87 s at $\beta = 1.0$.

\begin{lstlisting}[style=py,firstnumber=344,caption={\texttt{tesifl/strategies/sage\_smart.py}, lines 344--381}]
    def _sensor_soc(self, cid: int) -> float:
        """SoC che il device legge dalla propria batteria."""
        return float(self.world.snapshot(cid).soc)

    def _subscribe(self) -> None:
        """Iscrizione: dati statici e SoC iniziale di ogni client."""
        cids = list(self._node_to_cid.values())
        for c in cids:
            # il primo "report" e' l'iscrizione, al tempo 0, col workload medio
            # della popolazione (non entra nella media dei workload del client:
            # non e' una lettura vera)
            self._record(c, self._sensor_soc(c), self.w_enroll, False, 0.0,
                         measured=False)
            self._k_cal[c] = 1.0
        nominal = [
            training_time_s(self.world.profile(c), int(self.epochs),
                            self._n[c], int(self.batch_size))
            for c in cids if self._n.get(c)
        ]
        if nominal:
            self._deadline_s = self.deadline_mult * float(np.median(nominal))
        self._subscribed = True
        log(INFO, "%s: %s client iscritti | deadline %.1f s "
            "(%.2f x mediana nominale, E=%s B=%s) | max %s ondate | batteria %s",
            self.variant, len(cids), self._deadline_s, self.deadline_mult,
            self.epochs, self.batch_size, self.max_waves,
            BATTERY_NAMES[self.world.battery])

    def _record(self, cid: int, soc: float, w: float, chg: bool, t: float,
                measured: bool = True) -> None:
        """Un report del client: diventa anche lo stato noto."""
        if measured:
            self._w_sum[cid] = self._w_sum.get(cid, 0.0) + float(w)
            self._w_cnt[cid] = self._w_cnt.get(cid, 0) + 1
        self._soc_rep[cid], self._w_rep[cid] = float(soc), float(w)
        self._chg_rep[cid], self._t_rep[cid] = bool(chg), float(t)
        self._soc_known[cid], self._w_known[cid] = float(soc), float(w)
        self._chg_known[cid] = bool(chg)
\end{lstlisting}

\newpage
%======================================================================
\section{Step 1: forecast of the state}
\label{sec:forecast}

At the start of every selection the server computes, for every phone, the SoC
and the workload it expects now. The server clock $t^{\mathrm{now}}$ is the sum
of the durations of the previous rounds, which the server knows because it
waited for the replies. Let $s^{0}_i$, $w^{\mathrm{rep}}_i$, $c_i$ and
$t^{\mathrm{rep}}_i$ be the SoC, workload, charging flag and time of the last
report, $\Delta_i = t^{\mathrm{now}} - t^{\mathrm{rep}}_i$, and $\bar{w}_i$ the
mean of the workloads reported by the phone ($\bar{w}$ if none). If
$\Delta_i \le 0$ the forecast is the report itself. Otherwise the interval is
split into $J = \lceil \Delta_i / 300 \rceil$ steps of $\delta = \Delta_i / J$.

The workload relaxes from the last value towards the mean of the phone:
\begin{equation}
  w_j = \bar{w}_i + (w^{\mathrm{rep}}_i - \bar{w}_i)\, e^{-(j + 1/2)\,\delta / \tau},
  \qquad
  \hat{w}_i = \bar{w}_i + (w^{\mathrm{rep}}_i - \bar{w}_i)\, e^{-\Delta_i / \tau} .
\end{equation}

The SoC is the expected value over three branches with masses and SoC
$(m_{\mathrm{d}}, s_{\mathrm{d}})$ discharging, $(m_{\mathrm{c}},
s_{\mathrm{c}})$ charging, $(m_{\mathrm{f}}, s_{\mathrm{f}})$ charge completed
and discharging again. Initially $(m_{\mathrm{d}}, m_{\mathrm{c}}) = (1, 0)$
with $s_{\mathrm{d}} = s^0_i$, or $(0, 1)$ with $s_{\mathrm{c}} = s^0_i$ if the
phone was charging; $m_{\mathrm{f}} = 0$, $s_{\mathrm{f}} = 1$. At every step
$j$, in this order:
\begin{enumerate}
  \item the mass $\Delta m = m_{\mathrm{d}} \bigl(1 - e^{-\lambda (1 -
        s_{\mathrm{d}}) \delta / 3600}\bigr)$ moves from discharging to
        charging (the hourly plug-in model of the users, the same as
        Equation~\eqref{eq:plug}):
        $s_{\mathrm{c}} \leftarrow \frac{m_{\mathrm{c}} s_{\mathrm{c}} + \Delta m \, s_{\mathrm{d}}}{m_{\mathrm{c}} + \Delta m}$,
        $m_{\mathrm{c}} \leftarrow m_{\mathrm{c}} + \Delta m$,
        $m_{\mathrm{d}} \leftarrow m_{\mathrm{d}} - \Delta m$;
  \item the discharging branch pays the idle:
        $s_{\mathrm{d}} \leftarrow \max\bigl(0, s_{\mathrm{d}} + \Delta s(P^{\mathrm{idle}}(w_j), \delta, s_{\mathrm{d}})\bigr)$,
        with the battery of the phone; the same for $s_{\mathrm{f}}$ if
        $m_{\mathrm{f}} > 10^{-9}$;
  \item the charging branch gains $\gamma \delta / 3600$ (at most 1); when it
        reaches 1 its mass moves to the completed branch:
        $s_{\mathrm{f}} \leftarrow \frac{m_{\mathrm{f}} s_{\mathrm{f}} + m_{\mathrm{c}}}{m_{\mathrm{f}} + m_{\mathrm{c}}}$,
        $m_{\mathrm{f}} \leftarrow m_{\mathrm{f}} + m_{\mathrm{c}}$, $m_{\mathrm{c}} \leftarrow 0$.
\end{enumerate}
Finally $\hat{s}_i = \min\bigl(1, \max(0, m_{\mathrm{d}} s_{\mathrm{d}} +
m_{\mathrm{c}} s_{\mathrm{c}} + m_{\mathrm{f}} s_{\mathrm{f}})\bigr)$.
$\hat{s}_i$ is used for the ranking and for the energy limit of the envelope;
$\hat{w}_i$ is used for the deadline limit of the envelopes of the proposals.
The charging flag known to the server is the one of the last report.

Two columns of the CSV measure the error, without entering any decision:
\cf{smart\_info\_err}, the mean of $|\hat{s}_i - s_i|$ over $\mathcal{A}_r$, and
\cf{smart\_w\_err}, the mean of $|\hat{w}_i - w_i|$. With the ablation
\cf{noest} the known state is the last report as it is.

\begin{lstlisting}[style=py,firstnumber=25,caption={\texttt{tesifl/strategies/sage\_smart.py}, lines 25--27}]
# passo massimo con cui il server integra la previsione (passo 1): il mondo
# scarica l'idle una volta per round, e un round dura 100-300 s
_ESTIMATE_STEP_S = 300.0
\end{lstlisting}

\begin{lstlisting}[style=py,firstnumber=383,caption={\texttt{tesifl/strategies/sage\_smart.py}, lines 383--464}]
    def _forecast(self, cid: int) -> tuple[float, float]:
        """Passo 1: SoC e workload ATTESI adesso.

            w(t)   = w_med + (w_rep - w_med) * exp(-(t - t_rep) / tau)
            SoC(t) = valore atteso su tre rami, a passi di _ESTIMATE_STEP_S:
                     in scarica (idle con w(t)), in carica (C-rate fino al
                     100%), carica finita e di nuovo in scarica.

        w_med e' la media dei workload riportati dal client (prima del primo
        report, quella della popolazione). Dal ramo in scarica passa in carica,
        a ogni passo, la frazione p = 1 - exp(-lambda (1 - SoC) dt): e' il
        modello orario degli utenti (Ferreira 2011), un dato di POPOLAZIONE
        che il server conosce come conosce la batteria del device. Un client
        in carica all'ultimo report resta in carica fino al 100%. L'idle e'
        integrato a passi di al piu' _ESTIMATE_STEP_S, come lo scarica il mondo
        round per round: con la batteria da datasheet il calo dipende dal SoC.

        [B] PERCHE'. Una stima che scarica sempre e non ricarica mai fa
        scendere in classifica un client non contattato da ore anche se nel
        frattempo, in media, ha ricaricato (+2 punti l'ora a SoC 0.6). E un
        workload fermo all'ultimo valore sbaglia: dopo 5 round lo stato e' lo
        stesso solo nel 51% dei casi.
        """
        soc0 = float(self._soc_rep[cid])
        dt = max(0.0, self._clock_s - self._t_rep.get(cid, 0.0))
        w_rep = float(self._w_rep.get(cid, self.w_enroll))
        w_med = (self._w_sum[cid] / self._w_cnt[cid]) if self._w_cnt.get(cid) else self.w_enroll
        if dt <= 0.0:
            return soc0, w_rep
        prof = self.world.profile(cid)
        bat = self.world.battery
        recharge = bool(getattr(self.world, "_recharge_enabled", False))
        n = max(1, math.ceil(dt / _ESTIMATE_STEP_S))
        h = dt / n
        # masse e SoC dei tre rami: scarica (d), in carica (c), carica finita (f)
        if self._chg_rep.get(cid, False):
            m_d, s_d, m_c, s_c = 0.0, soc0, 1.0, soc0
        else:
            m_d, s_d, m_c, s_c = 1.0, soc0, 0.0, soc0
        m_f, s_f = 0.0, 1.0
        for i in range(n):
            w = w_med + (w_rep - w_med) * math.exp(-(i + 0.5) * h / self.load_tau_s)
            p_idle = idle_power_w(prof, w)
            if recharge and m_d > 0.0:
                moved = m_d * (1.0 - math.exp(-RECHARGE_RATE_PER_H * (1.0 - s_d) * h / 3600.0))
                if moved > 0.0:
                    s_c = (m_c * s_c + moved * s_d) / (m_c + moved)
                    m_c, m_d = m_c + moved, m_d - moved
            if m_d > 0.0:
                s_d = max(0.0, s_d + battery_delta_soc(bat, p_idle, h, prof, s_d))
            if m_f > 1e-9:
                s_f = max(0.0, s_f + battery_delta_soc(bat, p_idle, h, prof, s_f))
            if m_c > 0.0:
                s_c = min(1.0, s_c + CHARGE_C_RATE * h / 3600.0)
                if s_c >= 1.0:
                    s_f = (m_f * s_f + m_c) / (m_f + m_c)
                    m_f, m_c = m_f + m_c, 0.0
        soc = m_d * s_d + m_c * s_c + m_f * s_f
        w_now = w_med + (w_rep - w_med) * math.exp(-dt / self.load_tau_s)
        return float(min(1.0, max(0.0, soc))), float(w_now)

    def _refresh_state(self, available: list[int]) -> None:
        """Passo 1: lo stato con cui il server ordina i client."""
        for c in self._node_to_cid.values():
            if self.forecast:
                self._soc_known[c], self._w_known[c] = self._forecast(c)
            else:
                # ablazione noest: l'ultimo report cosi' com'e', anche se
                # vecchio di molti round
                self._soc_known[c] = self._soc_rep[c]
                self._w_known[c] = self._w_rep[c]
            self._chg_known[c] = self._chg_rep[c]
        # [B] verifica, non input: quanto lo stato noto si discosta da quello
        # che il device leggerebbe adesso, sul pool selezionabile
        pool = [self._node_to_cid[n] for n in available]
        if pool:
            self._info_err = float(np.mean(
                [abs(self._soc_known[c] - self._sensor_soc(c)) for c in pool]))
            if self._protocol_cols:
                self._w_err = float(np.mean(
                    [abs(self._w_known[c] - self.world.workload.utilization(c))
                     for c in pool]))
\end{lstlisting}

\newpage
%======================================================================
\section{Step 2: ranking}
\label{sec:ranking}

Every phone of $\mathcal{A}_r$ receives the score
\begin{equation}
  \sigma_i = a \, \hat{s}_i + b \, D_i + c \, \theta_i ,
  \qquad
  \theta_i = \min\Bigl(1, \frac{r - \ell_i}{\Theta}\Bigr) ,
\end{equation}
where $\ell_i$ is the last round in which phone $i$ was selected or refused a
proposal, and $\ell_i = 0$ if neither ever happened. The workload does not
enter the ranking.

The divergence comes from the class histograms that the phones report after
every training (counts over all their training samples). With $\pi_i$ the
normalized histogram, the Jensen-Shannon divergence in base 2, with additive
smoothing $\varepsilon = 10^{-9}$, is
\begin{equation}
  \tilde{\pi} = \frac{\pi + \varepsilon}{\sum (\pi + \varepsilon)}, \quad
  \bar{\pi} = \tfrac{1}{2}(\tilde{\pi}_i + \tilde{\pi}_j), \quad
  \mathrm{JS}(\pi_i, \pi_j) = \tfrac{1}{2}\mathrm{KL}_2(\tilde{\pi}_i \| \bar{\pi})
                            + \tfrac{1}{2}\mathrm{KL}_2(\tilde{\pi}_j \| \bar{\pi}) ,
\end{equation}
and $D_i = \frac{1}{H - 1} \sum_{j \ne i} \mathrm{JS}(\pi_i, \pi_j)$ over the $H$
phones with a known histogram ($D_i = 0$ for all if $H < 2$). When a histogram arrives, $D$ is recomputed for
all of them. A phone without a histogram gets the mean of the known $D_j$, or 0
if there is none, as for every phone at round 1.

\begin{lstlisting}[style=py,firstnumber=13,caption={\texttt{tesifl/strategies/sage.py}, lines 13--19}]
def _js(p: np.ndarray, q: np.ndarray) -> float:
    """Jensen-Shannon in base 2, con smoothing additivo come nel paper."""
    eps = 1e-9
    p, q = (p + eps) / (p + eps).sum(), (q + eps) / (q + eps).sum()
    m = 0.5 * (p + q)
    kl = lambda a, b: float(np.sum(a * (np.log(a) - np.log(b))) / np.log(2.0))
    return 0.5 * kl(p, m) + 0.5 * kl(q, m)
\end{lstlisting}

\begin{lstlisting}[style=py,firstnumber=40,caption={\texttt{tesifl/strategies/sage\_ablation.py}, lines 40--67}]
    def __init__(
        self,
        *args,
        sage_a: float = 0.5,
        sage_b: float = 0.5,
        sage_k: int = 2,
        half_epochs_threshold: float = 0.3,
        max_fraction: float = 0.2,   # k_per_round / N
        class_distributions: dict[int, np.ndarray] | None = None,
        **kwargs,
    ) -> None:
        super().__init__(*args, **kwargs)
        self.a, self.b = float(sage_a), float(sage_b)
        if abs(self.a + self.b - 1.0) > 1e-9:
            raise ValueError("a + b = 1")
        self.k = int(sage_k)
        self.half_th = float(half_epochs_threshold)
        self.max_frac = float(max_fraction)

        self._dist: dict[int, np.ndarray] = {}   # cid -> distribuzione classi
        self._div: dict[int, float] = {}         # cid -> JSD media
        self._half: set[int] = set()

        if class_distributions:
            for cid, d in class_distributions.items():
                arr = np.asarray(d, dtype=float)
                self._dist[int(cid)] = arr / arr.sum()
            self._recompute_div()
\end{lstlisting}

\begin{lstlisting}[style=py,firstnumber=74,caption={\texttt{tesifl/strategies/sage\_ablation.py}, lines 74--77}]
    def _divergence(self, cid: int) -> float:
        if cid in self._div:
            return self._div[cid]
        return float(np.mean(list(self._div.values()))) if self._div else 0.0
\end{lstlisting}

\begin{lstlisting}[style=py,firstnumber=140,caption={\texttt{tesifl/strategies/sage\_ablation.py}, lines 140--164}]
    def _on_round_result(
        self, cid: int, n_ex: int, dt_s: float, e_wh: float, e_lin_wh: float,
        reply: Message, server_round: int,
    ) -> None:
        # nessuna contabilita' energetica locale: il SoC lo tiene il world state
        dist = self._metric(reply, "class-dist", default=None)
        if dist is None:
            return
        arr = np.asarray(list(dist), dtype=float)
        if arr.sum() <= 0:
            return
        self._dist[cid] = arr / arr.sum()
        self._recompute_div()

    def _recompute_div(self) -> None:
        cids = sorted(self._dist)
        n = len(cids)
        if n < 2:
            self._div = {c: 0.0 for c in cids}
            return
        mat = np.zeros((n, n))
        for i in range(n):
            for j in range(i + 1, n):
                mat[i, j] = mat[j, i] = _js(self._dist[cids[i]], self._dist[cids[j]])
        self._div = {c: float(mat[i].sum() / (n - 1)) for i, c in enumerate(cids)}
\end{lstlisting}

\begin{lstlisting}[style=py,firstnumber=527,caption={\texttt{tesifl/strategies/sage\_smart.py}, lines 527--550}]
    def _resource(self, cid: int) -> float:
        """SoC usato nello score: quello NOTO al server, non quello vero."""
        return self._soc_known.get(cid, float(self.world.initial_soc(cid)))

    def _n_of(self, cid: int) -> int | None:
        if cid in self._n:
            return self._n[cid]
        return int(np.median(list(self._n.values()))) if self._n else None

    def _batch_for(self, cid: int) -> int:
        """B preferito: il piu' piccolo dell'envelope (passo 6 nella docstring
        della classe). Il B effettivo lo sceglie il client e lo dichiara."""
        return int(self.batch_choices[0])

    def _staleness(self, cid: int, server_round: int) -> float:
        """Round dall'ultima selezione, normalizzati su stale_max.

        Chi non e' mai stato scelto conta come scelto al round 0: la sua
        anzianita' cresce dal primo round come quella degli altri, e lo porta
        in classifica invece di escluderlo per sempre.
        """
        last = self._last_seen.get(cid)
        age = server_round if last is None else server_round - last
        return float(min(1.0, age / max(1, self.stale_max)))
\end{lstlisting}

\newpage
%======================================================================
\section{Step 3: number of slots}
\label{sec:slots}

The server solves
\begin{equation}
  \max_{y \in \{0,1\}^{m_r}} \sum_{i \in \mathcal{A}_r} y_i \, \sigma_i
  \quad \text{s.t.} \quad
  k \le \sum_i y_i \le h_r = \min\bigl(k + r - 1, \lfloor 0.2 \max(m_r, n^{\mathrm{alive}}_r) \rfloor, m_r\bigr)
\end{equation}
with PuLP and the CBC solver, $k = 2$. Without broadcast $n^{\mathrm{alive}}_r =
m_r$. The scores are strictly positive (the staleness term is at least
$c / \Theta$, because $\ell_i < r$), so the solution takes exactly $h_r$ phones, the first $h_r$
by score: 2 slots at round 1, 3 at round 2, 20 from round 19 while at least
100 phones are alive (19 with 95 to 99 alive, and so on). Only the number of
slots is used afterwards. If the problem is infeasible ($m_r < k$) the run
ends. The phones are then sorted by decreasing $\sigma_i$ (ties keep the order
of the phone index) and the proposals follow this ranking.

The argument \cf{k} of \cf{\_select\_clients}, $\min(20, m_r)$, is only checked
to be non-zero; the slots are $h_r$.

\begin{lstlisting}[style=py,firstnumber=553,caption={\texttt{tesifl/strategies/sage\_smart.py}, lines 553--643}]
    def _select_clients(
        self, available: list[int], k: int, server_round: int
    ) -> list[int]:
        if not self._subscribed:
            self._subscribe()
        self._reports = {}
        self._bound = {}
        self._queried = set()
        self._overhead_s = 0.0
        self._waves = 0
        self._n_queried = 0
        self._n_refused = 0
        self._n_deferred = 0
        self._n_noreply = 0
        self._extra_wait_s = 0.0
        if not k or not available:
            log(INFO, "%s round %s: pool vuoto", self.variant, server_round)
            return []
        self._refresh_state(available)
        cid_of = self._node_to_cid
        pool = available
        m = len(pool)
        # tetto sui VIVI, non sui soli raggiunti dal broadcast (vedi
        # PhysicalSAGEAblation._select_clients)
        hi = min(self.k + server_round - 1,
                 int(np.floor(self.max_frac * max(m, self._n_alive))), m)
        lo = self.k

        score = {
            n: self.a * self._resource(cid_of[n])
            + self.b * self._divergence(cid_of[n])
            + self.stale_w * self._staleness(cid_of[n], server_round)
            for n in pool
        }

        prob = pulp.LpProblem("SAGE_Smart_Select", pulp.LpMaximize)
        y = [pulp.LpVariable(f"y_{i}", cat="Binary") for i in range(m)]
        prob += pulp.lpSum(y[i] * score[pool[i]] for i in range(m))
        prob += pulp.lpSum(y) >= lo
        prob += pulp.lpSum(y) <= hi
        prob.solve(pulp.PULP_CBC_CMD(msg=0))

        selected = [pool[i] for i in range(m) if int(pulp.value(y[i]) or 0) == 1]
        if not selected:
            self._charge_control([])
            log(INFO, "%s round %s: ILP infeasible con pool %s (lo=%s, hi=%s): "
                "pool esaurito", self.variant, server_round, m, lo, hi)
            return []

        # [B] l'ILP con soli vincoli di cardinalita' sceglie i top per score:
        # le proposte scorrono la stessa classifica
        ranking = sorted(pool, key=lambda n: -score[n])
        if self.negotiation:
            selected = self._negotiate_pick(ranking, len(selected), server_round)
        else:
            # ablazione noneg: i primi della classifica senza chiedere, nessun
            # rimpiazzo e nessuno scambio di controllo (waves = 0, overhead 0)
            selected = ranking[:len(selected)]
            if self.proposal_nonresp:
                # [B] senza broadcast e senza proposte: chi non e'
                # raggiungibile non riceve il modello e il suo posto resta
                # vuoto (il round dura comunque fino alla deadline)
                lost = {n for n in selected
                        if self._bcast_rng.random() < self.bcast_nonresp}
                self._n_noreply += len(lost)
                selected = [n for n in selected if n not in lost]
        if not selected:
            self._charge_control([])
            log(INFO, "%s round %s: nessun candidato ha accettato in %s "
                "ondate: pool esaurito", self.variant, server_round,
                self._waves)
            return []

        self._charge_control(selected)
        self._overhead_s = self._waves * ctrl_latency_s() + self._extra_wait_s
        for n in selected:
            self._last_seen[cid_of[n]] = server_round
        # envelope con lo stato appena riportato dai client che hanno accettato
        self._plan = {cid_of[n]: self._assign_work(cid_of[n]) for n in selected}

        E = [env.e_hi for env in self._plan.values()]
        never = sum(1 for c in cid_of.values() if c not in self._last_seen)
        log(INFO, "%s round %s: pool %s, %s selezionati (cap %s) | "
            "E_hi %s-%s (medio %.1f) | mai visti %s | interrogati %s, "
            "rifiuti %s, ondate %s%s",
            self.variant, server_round, m, len(selected), hi, min(E), max(E),
            float(np.mean(E)), never, self._n_queried, self._n_refused,
            self._waves,
            (f" | rinviati {self._n_deferred}, non raggiunti {self._n_noreply}, "
             f"errore w {self._w_err:.3f}") if self._protocol_cols else "")
        return selected
\end{lstlisting}

\newpage
%======================================================================
\section{Step 4: proposals (negotiation)}
\label{sec:proposals}

The server does not send the model to the first $h_r$ phones of the ranking.
It sends a proposal, the envelope of Section~\ref{sec:envelope} without the
model, as a Flower \cf{QUERY} message.
\begin{enumerate}
  \item Wave 1 goes to the first $\lceil \kappa h_r \rceil$ phones of the
        ranking (30 when $h_r = 20$). Their envelopes are computed with the
        forecast state.
  \item Every proposed phone is unreachable with probability $p_{\mathrm{nr}} =
        0.10$ (own random generator, seed of the run + 104729). It receives
        nothing, pays nothing, does not answer and does not count as
        contacted. If at least one phone of the wave is unreachable, the wave
        waits $t_{\mathrm{to}} - t_{\mathrm{ctrl}} = 0.665$ s more than
        $t_{\mathrm{ctrl}}$.
  \item A phone that answers sends its current SoC, workload and charging
        flag, which become its last report (time $t^{\mathrm{now}}$) and its
        known state, and accepts or refuses. It refuses only for energy: it
        accepts if it is charging, or if for at least one $B \in \mathcal{B}$
        the work $E_{\mathrm{lo}} = 2$ epochs plus the probe of 256 samples,
        followed by the model upload, leaves its predicted SoC at or above
        $s_{\min}$ (Equations~\eqref{eq:ttrain} and \eqref{eq:predict}, with its
        current workload).
  \item A refusal sets $\ell_i = r$: the staleness of the phone restarts from
        zero.
  \item If the accepted phones are fewer than $h_r$, wave 2 goes to the next
        $\lceil \kappa (h_r - \text{accepted}) \rceil$ phones of the ranking
        not yet asked. There are at most $W = 2$ waves. If after them the
        accepted phones are still fewer than $h_r$, the round goes on with
        fewer phones; if none accepted, the pool is exhausted and the run
        ends.
\end{enumerate}

The energy of the control messages is charged after the fresh choice: a phone
that is selected pays $1.97 \cdot 10^{-3}$ J (the exchange shares the radio
session of the model download), every other phone that answered pays 2.138 J.
The control overhead of the round, which delays the training of the selected
phones and is paid as idle by all the others, is
\begin{equation}
  t^{\mathrm{ovh}}_r = (\text{waves}) \cdot t_{\mathrm{ctrl}}
                     + (\text{waves with an unreachable phone}) \cdot (t_{\mathrm{to}} - t_{\mathrm{ctrl}}) .
\end{equation}

\begin{lstlisting}[style=py,firstnumber=658,caption={\texttt{tesifl/strategies/sage\_smart.py}, lines 658--697}]
    def _negotiate_pick(self, ranking: list[int], slots: int,
                        server_round: int) -> list[int]:
        """Passi 4-5: rosa piu' larga dei posti, scelta sullo stato fresco.

        Ondata 1: proposte ai primi ceil(shortlist * slots) della classifica.
        Chi accetta entra nella rosa; chi rifiuta (solo per energia) perde
        l'anzianita'; chi non risponde non e' raggiungibile e non conta come
        contatto. Se i posti non sono coperti, ondata 2 coi successivi, ancora
        shortlist volte quelli che mancano. Poi il server sceglie i `slots`
        migliori per _fresh_score: gli altri della rosa sono RINVIATI, tengono
        l'anzianita' e restano in cima al round dopo.
        """
        cid_of = self._node_to_cid
        accepted: list[int] = []
        asked: set[int] = set()
        need = slots
        while need > 0 and self._waves < self.max_waves:
            batch = [n for n in ranking if n not in asked][: int(math.ceil(self.shortlist * need))]
            if not batch:
                break
            self._waves += 1
            asked.update(batch)
            envs = {cid_of[n]: self._assign_work(cid_of[n]) for n in batch}
            replies = self._propose(batch, envs)
            for n in batch:
                r = replies.get(n)
                if r is None:
                    continue
                if int(self._rep(r, "smart-accept", 0)) == 1:
                    accepted.append(n)
                else:
                    self._n_refused += 1
                    # [B] contattato: l'anzianita' riparte da zero, cosi' non
                    # torna in testa alla classifica al round successivo
                    self._last_seen[cid_of[n]] = server_round
            need = slots - len(accepted)
        picked = sorted(accepted,
                        key=lambda n: -self._fresh_score(cid_of[n], server_round))[:slots]
        self._n_deferred = len(accepted) - len(picked)
        return picked
\end{lstlisting}

\begin{lstlisting}[style=py,firstnumber=466,caption={\texttt{tesifl/strategies/sage\_smart.py}, lines 466--510}]
    def _propose(self, nodes: list[int],
                 envs: dict[int, Envelope]) -> dict[int, Message]:
        """Passo 4: proposte (MessageType.QUERY) e lettura delle risposte.

        Le risposte aggiornano subito lo stato noto. Il costo energetico NON
        si addebita qui: dipende da chi verra' selezionato (sessione radio
        condivisa col download del modello o no), vedi _charge_control.
        """
        if self.proposal_nonresp:
            # [B] senza broadcast nessuno ha verificato prima chi e'
            # raggiungibile. La frazione broadcast_nonresp non riceve la
            # proposta (fuori copertura, spento): non paga nulla, non risponde,
            # e l'ondata si chiude al timeout.
            lost = {n for n in nodes if self._bcast_rng.random() < self.bcast_nonresp}
            if lost:
                self._n_noreply += len(lost)
                self._extra_wait_s += max(0.0, self.bcast_timeout_s - ctrl_latency_s())
                nodes = [n for n in nodes if n not in lost]
        msgs: list[Message] = []
        for n in nodes:
            cid = self._node_to_cid[n]
            cfg = ConfigRecord({"cid": int(cid)})
            for key, val in envs[cid].to_config().items():
                cfg[key] = val
            rec = RecordDict({self.configrecord_key: cfg,
                              "device": self._device_record(cid)})
            msgs.extend(self._construct_messages(rec, [int(n)], MessageType.QUERY))
        self._queried.update(nodes)
        self._n_queried += len(nodes)
        if self._grid is None or not msgs:
            return {}
        out: dict[int, Message] = {}
        for r in self._grid.send_and_receive(msgs, timeout=self.query_timeout_s):
            if not r.has_content():
                continue
            n = int(r.metadata.src_node_id)
            cid = self._node_to_cid.get(n)
            if cid is None:
                continue
            self._record(cid, float(self._rep(r, "smart-soc", self._soc_known[cid])),
                         float(self._rep(r, "smart-workload", self.w_busy)),
                         bool(int(self._rep(r, "smart-charging", 0))),
                         self._clock_s)
            out[n] = r
        return out
\end{lstlisting}

\begin{lstlisting}[style=py,firstnumber=512,caption={\texttt{tesifl/strategies/sage\_smart.py}, lines 512--525}]
    def _charge_control(self, selected: list[int]) -> None:
        """Addebita al mondo gli scambi della negoziazione del round.

        Chi e' stato selezionato riceve subito dopo il modello: lo scambio
        cade nella stessa sessione radio (merged). Chi non lo e' paga una
        sessione intera.
        """
        sel = set(selected)
        for n in self._queried:
            # [B] col broadcast attivo ogni candidato ha appena risposto alla
            # richiesta di stato: la sessione radio e' gia' aperta e la
            # proposta paga il solo payload, anche se viene rifiutata
            self.world.apply_control(self._node_to_cid[n],
                                     merged=(n in sel) or self.broadcast)
\end{lstlisting}

The phone answers with the handler \cf{query} of the client. Its reading is the
one attached by the server, taken before \cf{begin\_round}; the probe size in
the check is the one of the envelope (256), not capped by $n_i$ as it is at
training time.

\begin{lstlisting}[style=py,firstnumber=279,caption={\texttt{tesifl/tesiFL/client\_app.py}, lines 279--305}]
@app.query()
def query(msg: Message, context: Context):
    """Proposta di lavoro di sage_smart, prima del training (Fase 2b).

    Il server propone un envelope di (E, B): il client accetta se riesce a
    fare almeno il minimo senza intaccare la riserva, altrimenti rifiuta, e
    il server chiama un altro client al suo posto. La risposta porta anche
    SoC, workload e stato di carica letti adesso: il server li usa al posto
    della propria stima.

    Niente modello, niente dati: e' un messaggio di pochi byte. Il costo
    energetico (una sessione radio) lo addebita il mondo fisico.
    """
    cfg = msg.content["config"]
    dev = msg.content[DEVICE_KEY]
    reading = DeviceReading.from_record(dev)
    # [B] la taglia del dataset il client la conosce; nel simulatore arriva
    # nel record del device per non caricare i dati in una query
    n = int(dev["dataset-size"])
    metrics = {
        "smart-soc": float(reading.soc),
        "smart-workload": float(reading.workload),
        "smart-charging": int(reading.charging),
        "smart-accept": int(accepts_proposal(reading, Envelope.from_config(cfg), n)),
    }
    content = RecordDict({"metrics": MetricRecord(metrics)})
    return Message(content=content, reply_to=msg)
\end{lstlisting}

\begin{lstlisting}[style=py,firstnumber=128,caption={\texttt{tesifl/tesiFL/client\_policy.py}, lines 128--170}]
def statistical_utility(loss: float, n_classes: int = 10) -> float:
    if not math.isfinite(loss):
        return 1.0
    return float(min(1.0, max(0.0, loss / math.log(n_classes))))


def headroom(reading: DeviceReading, soc_min: float) -> float:
    if reading.charging:
        return 1.0
    return float(min(1.0, max(0.0, (reading.soc - soc_min) / max(1e-9, 1.0 - soc_min))))


def _energy_ok(reading: DeviceReading, soc_end: float, soc_min: float) -> bool:
    return reading.charging or soc_end >= soc_min


def can_afford_minimum(reading: DeviceReading, env: Envelope, n_samples: int) -> bool:
    """Il round minimo (E_lo, batch migliore, stima della loss inclusa) lascia
    la riserva intatta? Si verifica PRIMA della stima della loss: se la
    risposta e' no, non vale la pena spendere nemmeno quella."""
    if reading.charging:
        return True
    for b in env.batches:
        _, soc_end = reading.predict(env.e_lo, b, n_samples, env.probe_samples)
        if soc_end >= env.soc_min:
            return True
    return False


def accepts_proposal(reading: DeviceReading, env: Envelope, n_samples: int) -> bool:
    """Controllo di accettazione della proposta (passo 4 di sage_smart).

    Il client accetta se riesce a fare almeno il lavoro minimo dell'envelope
    senza intaccare la riserva, con lo stato di ADESSO; in carica accetta
    sempre. E' lo stesso controllo che fara' di nuovo al momento del
    training: qui serve a rifiutare PRIMA di scaricare il modello, cosi' il
    server puo' chiamare un altro client nello stesso round.

    [B] rifiuta SOLO per energia. Chi non sta nella deadline accetta e fa i
    passi che ci stanno (lavoro parziale, vedi decide): il rifiuto per
    deadline (05-06/11) escludeva i telefoni lenti con molti dati.
    """
    return can_afford_minimum(reading, env, n_samples)
\end{lstlisting}

\newpage
%======================================================================
\section{Step 5: fresh choice}
\label{sec:fresh}

Among the phones that accepted, the server keeps the $h_r$ with the highest
fresh score, computed with the state they have just sent:
\begin{equation}
  \phi_i = \bigl(a \, s'_i + b \, D_i + c \, \theta_i\bigr) \, (1 - w'_i),
  \qquad s'_i = \begin{cases} 1 & \text{if charging} \\ \text{reported SoC} & \text{otherwise} \end{cases}
\end{equation}
with $w'_i$ the reported workload. A charging phone counts as full because its
training does not touch the battery; a busy phone counts less because within
the deadline it would do fewer steps. The other phones that accepted are
deferred: $\ell_i$ does not change, so they keep their staleness, and their
known state is the fresh report. The selected phones get $\ell_i = r$, then
their envelopes are recomputed with the fresh state (Section~\ref{sec:envelope}).

\begin{lstlisting}[style=py,firstnumber=645,caption={\texttt{tesifl/strategies/sage\_smart.py}, lines 645--656}]
    def _fresh_score(self, cid: int, server_round: int) -> float:
        """Passo 5: lo score con lo stato APPENA riportato dal client.

        Stessa forma della classifica, ma con SoC, carica e workload della
        risposta alla proposta: un telefono in carica conta come pieno (la
        batteria non paga), e il punteggio e' pesato per 1 - w, il lavoro
        utile che puo' fare adesso.
        """
        soc = 1.0 if self._chg_rep.get(cid, False) else float(self._soc_rep[cid])
        base = (self.a * soc + self.b * self._divergence(cid)
                + self.stale_w * self._staleness(cid, server_round))
        return base * (1.0 - float(self._w_rep.get(cid, self.w_enroll)))
\end{lstlisting}

\newpage
%======================================================================
\section{Step 6: envelope}
\label{sec:envelope}

For phone $i$, with throughput reduced by the known workload $\hat{w}_i$
(forecast for the proposals, fresh report for the selected phones):
\begin{align}
  E^{\mathrm{dl}}_i &= \Bigl\lfloor \frac{T}{\min_{B \in \mathcal{B}} t^{\mathrm{tr}}_i(1, n_i, B)} \Bigr\rfloor ,\\
  d^{\mathrm{ep}}_i &= k_i \, \frac{(P^{\mathrm{tr}}_i / \eta) \, t^{\mathrm{tr}}_i(1, n_i, 16) / 3600}{V_0 Q_i}, \qquad
  d^{\mathrm{up}}_i = k_i \, \frac{(P^{\mathrm{comm}} / \eta) \, t^{\mathrm{comm}} / 3600}{V_0 Q_i} ,\\
  E^{\mathrm{en}}_i &= \Bigl\lfloor \frac{\hat{s}_i - s_{\min} - d^{\mathrm{up}}_i}{d^{\mathrm{ep}}_i} \Bigr\rfloor
      \quad (E^{\mathrm{en}}_i = E_{\mathrm{cap}} \text{ if the phone is charging}) ,\\
  E_{\mathrm{hi},i} &= \min\bigl(E_{\mathrm{cap}}, \max(E_{\mathrm{lo}}, \min(E_{\mathrm{cap}}, E^{\mathrm{dl}}_i, E^{\mathrm{en}}_i))\bigr) .
\end{align}
The energy limit is linear (Ideal battery) corrected by $k_i$, at the most
expensive batch per epoch (16); the deadline limit uses the fastest batch. If
$n_i$ is unknown the median of the known sizes is used. With $n = 400$ and
$w = \bar{w}$, $E^{\mathrm{dl}} = 2$, 4 and 5 for the three tiers, while
$E^{\mathrm{en}}$ is above 200 at SoC 0.75 (one epoch costs 0.19--0.26 points of
SoC, the model exchange 0.27--0.28): in practice the deadline limits $E_{\mathrm{hi}}$.
The CSV records which limit was active (\cf{cap}, \cf{deadline},
\cf{energy}).

The envelope contains $E_{\mathrm{lo}} = 2$, $E_{\mathrm{hi},i}$, $\mathcal{B}
= (16, 32, 64, 128)$, $T$, $s_{\min} = 0.20$, 256 probe samples and the flag of
the partial work. The training message also carries \cf{epochs}
$= E_{\mathrm{hi},i}$ and \cf{batch-size} $= 16$, which the client does not use
to train, and the sensor record of the phone, read by the world after
\cf{begin\_round}.

\begin{lstlisting}[style=py,firstnumber=705,caption={\texttt{tesifl/strategies/sage\_smart.py}, lines 705--774}]
    def _batches_for(self, cid: int) -> tuple[int, ...]:
        """Taglie di batch dell'envelope, dalla piu' piccola (preferita)."""
        return self.batch_choices

    def _assign_work(self, cid: int) -> Envelope:
        prof = self.world.profile(cid)          # statico: dal round 0
        batches = self._batches_for(cid)
        b_ref = self._batch_for(cid)
        peff = effective_profile(prof, self._w_known.get(cid, self.w_enroll))
        e_lo, e_cap = int(self.epochs_min), int(self.epochs_cap)
        n = self._n_of(cid)

        # vincolo di deadline: col batch piu' veloce dell'envelope
        e_dl = e_cap
        if n and math.isfinite(self._deadline_s):
            t_ep = min(training_time_s(peff, 1, n, b) for b in batches)
            e_dl = int(self._deadline_s // max(t_ep, 1e-9))

        # vincolo di energia: stima lineare corretta da k_i
        e_en = e_cap
        if n and not self._chg_known.get(cid, False):
            cap_wh = prof.battery_capacity_mah / 1000.0 * V_NOMINAL
            k = self._k_cal.get(cid, 1.0)
            d_ep = k * energy_wh_linear(
                training_power_w(prof), training_time_s(peff, 1, n, b_ref)
            ) / cap_wh
            d_up = k * energy_wh_linear(comm_power_w(), comm_time_s()) / cap_wh
            spare = self._resource(cid) - self.soc_min - d_up
            e_en = int(spare // d_ep) if d_ep > 0 else e_cap

        e_hi = int(np.clip(min(e_cap, e_dl, e_en), e_lo, e_cap))
        if e_hi >= e_cap:
            self._bound[cid] = "cap"
        elif e_dl <= e_en:
            self._bound[cid] = "deadline"
        else:
            self._bound[cid] = "energy"
        return Envelope(
            e_lo=e_lo, e_hi=e_hi, batches=batches,
            deadline_s=float(self._deadline_s) if math.isfinite(self._deadline_s)
            else 1e12,
            soc_min=float(self.soc_min), probe_samples=self.probe_samples,
            partial_work=self.partial_work,
        )

    def _per_client_config(
        self, cid: int, cfg: ConfigRecord, server_round: int
    ) -> None:
        env = self._plan.get(cid)
        if env is None:
            return
        for key, val in env.to_config().items():
            cfg[key] = val
        # [B] epochs e batch-size restano nel config come TETTO e batch di
        # riferimento: il client di sage_smart non li usa per allenarsi.
        cfg["epochs"] = int(env.e_hi)
        cfg["batch-size"] = int(self._batch_for(cid))

    def _device_record(self, cid: int) -> ConfigRecord | None:
        """Sensori del device: letti dal mondo PER CONTO del client.

        [B] nel messaggio di training e' chiamato dopo WorldState.begin_round:
        il client vede il workload e lo stato di carica di QUESTO round, il
        server nella selezione aveva visto quelli del precedente.
        [B] "dataset-size": il client conosce la taglia del proprio dataset;
        nel simulatore arriva qui per non dover caricare i dati in una query.
        """
        rec = DeviceReading.from_world(self.world, cid).to_record()
        rec["dataset-size"] = int(self._n_of(cid) or 0)
        return ConfigRecord(rec)
\end{lstlisting}

\begin{lstlisting}[style=py,firstnumber=54,caption={\texttt{tesifl/tesiFL/client\_policy.py}, lines 54--61}]
# chiavi del ConfigRecord con cui viaggia l'envelope
K_E_LO = "smart-e-lo"
K_E_HI = "smart-e-hi"
K_BATCHES = "smart-batches"
K_DEADLINE = "smart-deadline-s"
K_SOC_MIN = "smart-soc-min"
K_PROBE = "smart-probe-samples"
K_PARTIAL = "smart-partial-work"
\end{lstlisting}

\begin{lstlisting}[style=py,firstnumber=64,caption={\texttt{tesifl/tesiFL/client\_policy.py}, lines 64--111}]
@dataclass(frozen=True)
class Envelope:
    """Proposta del server: dentro questi limiti decide il client.

    partial_work: al training, se nemmeno E_lo sta nella deadline il client
    fa i passi di SGD che ci stanno (meno di E_lo epoche) invece di lavorare
    in ritardo (vedi decide). Viaggia nel messaggio solo se acceso.
    [B] 16/11: tolti refuse_deadline e refuse_load_only (rifiuto della
    proposta per la deadline, SAGE-smart plus/plus2-4): il client rifiuta
    solo per energia.
    """

    e_lo: int
    e_hi: int
    batches: tuple[int, ...]
    deadline_s: float
    soc_min: float
    probe_samples: int
    partial_work: bool = False

    def to_config(self) -> dict:
        cfg = {
            K_E_LO: int(self.e_lo),
            K_E_HI: int(self.e_hi),
            K_BATCHES: [int(b) for b in self.batches],
            K_DEADLINE: float(self.deadline_s),
            K_SOC_MIN: float(self.soc_min),
            K_PROBE: int(self.probe_samples),
        }
        if self.partial_work:
            cfg[K_PARTIAL] = 1
        return cfg

    @classmethod
    def from_config(cls, cfg) -> "Envelope":
        return cls(
            e_lo=int(cfg[K_E_LO]),
            e_hi=int(cfg[K_E_HI]),
            batches=tuple(int(b) for b in cfg[K_BATCHES]),
            deadline_s=float(cfg[K_DEADLINE]),
            soc_min=float(cfg[K_SOC_MIN]),
            probe_samples=int(cfg[K_PROBE]),
            partial_work=bool(int(cfg[K_PARTIAL])) if K_PARTIAL in cfg else False,
        )

    @staticmethod
    def present_in(cfg) -> bool:
        return K_E_LO in cfg
\end{lstlisting}

\newpage
%======================================================================
\section{Step 6, on the phone: decision and training}
\label{sec:client}

The client recognizes the envelope in the configuration and runs
\cf{\_train\_client\_driven}. The steps are:
\begin{enumerate}
  \item Reading of the sensors of this round: SoC, workload, charging flag
        (late binding). The local data are loaded; $n_i$ is the size of the
        training part; the probe size is $n_{\mathrm{p}} = \min(256, n_i)$.
  \item If not even $E_{\mathrm{lo}}$ is affordable (same check as the
        acceptance), the phone withdraws without measuring the loss: it pays
        the download and an isolated control message.
  \item Probe: mean cross-entropy $L_i$ of the global model on $n_{\mathrm{p}}$
        random local samples (generator seeded with $\text{seed} + 1000 r + i$,
        evaluation mode, no gradient, batches of 128).
  \item Decision (\cf{decide}):
        \begin{align}
          u_i &= \min\Bigl(1, \max\bigl(0, L_i / \ln |\mathcal{Y}|\bigr)\Bigr)
               \quad (\ln 100 = 4.605;\; u_i = 1 \text{ if } L_i \text{ is not finite}), \\
          h_i &= \begin{cases} 1 & \text{if charging} \\
                 \min\bigl(1, \max(0, (s_i - s_{\min}) / (1 - s_{\min}))\bigr) & \text{otherwise} \end{cases} \\
          E^{*}_i &= E_{\mathrm{lo}} + \operatorname{round}\bigl((E_{\mathrm{hi},i} - E_{\mathrm{lo}})\, u_i \, h_i\bigr)
        \end{align}
        (Python \cf{round}, halves to even). For $B$ in increasing order, and
        for $E$ from $E^{*}_i$ down to $E_{\mathrm{lo}}$, the first pair whose
        predicted time is within $T$ and whose predicted final SoC is at least
        $s_{\min}$ (or the phone is charging) is chosen. The prediction is
        \begin{equation}
          t = t^{\mathrm{tr}}_i(E, n_i, B) + t^{\mathrm{probe}}_i(n_{\mathrm{p}}), \quad
          s^{\mathrm{end}} = s_i + \Delta s(P^{\mathrm{tr}}_i, t, s_i) + \Delta s(P^{\mathrm{comm}}, t^{\mathrm{comm}}, \cdot)
          \label{eq:predict}
        \end{equation}
        with the workload of now and the battery of the phone, clamped to
        $[0, 1]$; a phone whose training would reach 0 is predicted at 0, and a
        charging phone keeps its SoC.
  \item Partial work: if no pair $(E, B)$ satisfies both conditions, but
        $E_{\mathrm{lo}}$ is affordable for the battery at some $B$,
        the phone takes the smallest $B$ at which the battery allows the steps
        that fit the deadline: the largest $\nu \le E_{\mathrm{lo}} \lceil
        n_i / B \rceil$ with predicted time within $T$, found by bisection, at
        least 1. It reports $E = E_{\mathrm{lo}}$ and a cap of $\nu$ steps.
  \item Otherwise, if $E_{\mathrm{lo}}$ is affordable, $E_{\mathrm{lo}}$ with
        the fastest affordable $B$, out of the deadline. Otherwise the phone
        withdraws (reserve).
  \item Training: SGD with learning rate 0.01 and momentum 0.9, shuffled
        batches of size $B$, for $E$ epochs, at most the planned steps $E
        \lceil n_i/B \rceil$ (or the cap of the partial work), and at most the
        steps that keep the predicted SoC at or above $s_{\min}$
        (\cf{affordable\_steps}, bisection; no limit if charging). The
        training can stop in the middle of an epoch.
  \item Reply: the model and the report. \cf{num-examples} $= n_i$;
        \cf{class-dist}, the class histogram; \cf{smart-epochs} $= E$;
        \cf{smart-batch} $= B$; \cf{smart-steps} $= \nu_i$, the steps really
        executed; the SoC at the start and the predicted SoC at the end for
        the work really done (time with $\nu_i / \lceil n_i/B \rceil$ epochs,
        plus the upload); the workload and the charging flag; the loss before
        and after; whether the deadline was met; with the partial work on,
        \cf{smart-partial} and \cf{smart-weight} $= \min(n_i, \nu_i B)$, the
        samples really seen.
\end{enumerate}

\begin{lstlisting}[style=py,firstnumber=68,caption={\texttt{tesifl/tesiFL/client\_app.py}, lines 68--72}]
    # sage_smart: il server ha mandato una PROPOSTA, non un'assegnazione
    if Envelope.present_in(config):
        return _train_client_driven(msg, model, device, config, partition_id,
                                    num_partitions, beta, seed, proximal_mu,
                                    dataset_name=dataset, augment=augment)
\end{lstlisting}

\begin{lstlisting}[style=py,firstnumber=103,caption={\texttt{tesifl/tesiFL/client\_app.py}, lines 103--138}]
def _class_counts(loader, n_classes: int = 10) -> list[int]:
    """Distribuzione delle classi locali, per la divergenza JS di SAGE.

    [B] su TUTTI i campioni: nel codice originale il bincount finale usava la
    sola variabile del ciclo, cioe' l'ultimo batch.
    """
    counts = [0] * n_classes
    for _b in loader:
        for _lbl in _b["label"].tolist():
            if 0 <= int(_lbl) < n_classes:
                counts[int(_lbl)] += 1
    return counts


def _probe_loss(model, dataset, n_probe: int, device, generator) -> float:
    """Loss del modello GLOBALE su n_probe campioni locali estratti a caso.

    Fase 5, passo 2 della proposta: alta loss = i dati di questo client sono
    ancora mal rappresentati nel modello globale = vale la pena lavorare.
    [B] in eval() e senza gradiente: in train() il forward aggiornerebbe le
    statistiche di BatchNorm, e l'update inviato al server ne porterebbe
    traccia anche se il client poi non si allenasse.
    """
    idx = torch.randperm(len(dataset), generator=generator)[:n_probe].tolist()
    loader = DataLoader(Subset(dataset, idx), batch_size=128)
    criterion = torch.nn.CrossEntropyLoss(reduction="sum")
    model.eval()
    tot, n = 0.0, 0
    with torch.no_grad():
        for batch in loader:
            images = batch["img"].to(device)
            labels = batch["label"].to(device)
            tot += float(criterion(model(images), labels).item())
            n += int(labels.numel())
    model.train()
    return tot / max(1, n)
\end{lstlisting}

\begin{lstlisting}[style=py,firstnumber=141,caption={\texttt{tesifl/tesiFL/client\_app.py}, lines 141--276}]
def _train_client_driven(msg: Message, model, device, config, partition_id: int,
                         num_partitions: int, beta: float, seed: int,
                         proximal_mu: float, dataset_name: str = "cifar10",
                         augment: bool = False) -> Message:
    """sage_smart: il client decide (E, B) dentro l'envelope del server.

    Fasi 5-7 della proposta:
      5. rilegge il proprio stato (late binding), stima la loss locale,
         sceglie (E, B) oppure si ritira;
      6. si allena con una guardia sul budget energetico;
      7. risponde con l'update e con cio' che ha fatto davvero: E e B scelti,
         passi eseguiti, SoC iniziale e finale, loss prima e dopo.

    Il SoC finale e' quello che il fuel gauge del device prevede per il
    lavoro svolto (training + upload): vedi device/fuel_gauge.py.
    """
    if DEVICE_KEY not in msg.content:
        raise KeyError(
            "sage_smart: il messaggio non contiene i sensori del device "
            f"(record '{DEVICE_KEY}'): la strategia deve allegarli"
        )
    # [B] LATE BINDING: la lettura e' quella di adesso, non quella che il
    # server aveva al momento della selezione
    reading = DeviceReading.from_record(msg.content[DEVICE_KEY])
    env = Envelope.from_config(config)
    server_round = int(config["server-round"])

    n_cls = num_classes(dataset_name)
    # [B] con augment la stima della loss (5b) vede immagini aumentate come
    # il training: misura quanto resta da imparare sull'obiettivo vero, che
    # con l'augmentation non va a zero per memorizzazione
    trainloader, _ = load_data(partition_id, num_partitions, env.batches[0],
                               beta, seed, dataset=dataset_name, augment=augment)
    dataset = trainloader.dataset
    n = len(dataset)
    env = dataclasses.replace(env, probe_samples=min(env.probe_samples, n))

    base = {
        "smart-soc-start": float(reading.soc),
        "smart-workload": float(reading.workload),
        "smart-charging": int(reading.charging),
    }

    def optout(probe: int, loss_pre: float, e_target: int) -> Message:
        """Ritiro dal round: nessun update, la riserva resta intatta.

        [B] num-examples = 0: il server lo esclude dalla media, ma addebita
        comunque la comunicazione (il modello e' stato scaricato) e, se
        c'e' stata, la stima della loss.
        """
        b = env.batches[0]
        dt = reading.work_time_s(0, b, n, probe)
        _, soc_final = reading.soc_after(dt, b)
        metrics = {
            **base,
            "num-examples": 0,
            "train_loss": 0.0,
            "smart-optout": 1,
            "smart-epochs": 0,
            "smart-batch": int(b),
            "smart-steps": 0,
            "smart-probe": int(probe),
            "smart-soc-final": float(soc_final),
            "smart-e-target": int(e_target),
        }
        if loss_pre == loss_pre:          # non NaN
            metrics["smart-loss-pre"] = float(loss_pre)
        content = RecordDict({"metrics": MetricRecord(metrics)})
        return Message(content=content, reply_to=msg)

    # 5a. il minimo e' sostenibile? Se no, inutile anche stimare la loss
    if not can_afford_minimum(reading, env, n):
        return optout(probe=0, loss_pre=float("nan"), e_target=0)

    # 5b. utilita' statistica: loss del modello globale sui dati locali
    model.to(device)
    gen = torch.Generator().manual_seed(seed + 1000 * server_round + partition_id)
    loss_pre = _probe_loss(model, dataset, env.probe_samples, device, gen)

    # 5c. decisione
    plan = decide(reading, env, n, loss_pre, n_classes=n_cls)
    if plan.optout:
        return optout(probe=env.probe_samples, loss_pre=loss_pre,
                      e_target=plan.e_target)

    # 6. training con guardia sul budget energetico
    loader = DataLoader(dataset, batch_size=plan.batch, shuffle=True)
    spe = len(loader)                        # ceil(n / B): passi per epoca
    planned = plan.epochs * spe
    # [B] 06/11 lavoro parziale (env.partial_work): al massimo i passi che
    # stanno nella deadline, vedi client_policy.decide
    if plan.steps_cap > 0:
        planned = min(planned, plan.steps_cap)
    budget = affordable_steps(reading, env, n, plan.batch, cap=planned)
    stats: dict = {}
    train_loss, loss_sq_sum, loss_n = train_fn(
        model, loader, plan.epochs, config["lr"], device,
        proximal_mu=proximal_mu, max_steps=budget, stats=stats,
    )
    steps = int(stats["steps"])

    # 7. SoC finale previsto dal fuel gauge per il lavoro EFFETTIVO
    dt = reading.work_time_s(steps / spe, plan.batch, n, env.probe_samples)
    _, soc_final = reading.soc_after(dt, plan.batch)

    metrics = {
        **base,
        "train_loss": train_loss,
        "num-examples": n,
        "class-dist": _class_counts(loader, n_cls),
        "loss-sq-sum": loss_sq_sum,
        "loss-n": loss_n,
        "smart-optout": 0,
        "smart-epochs": int(plan.epochs),
        "smart-batch": int(plan.batch),
        "smart-steps": steps,
        "smart-steps-budget": int(budget),
        "smart-budget-hit": int(steps < planned),
        "smart-probe": int(env.probe_samples),
        "smart-soc-final": float(soc_final),
        "smart-time-s": float(dt),
        "smart-deadline-ok": int(dt <= env.deadline_s),
        "smart-e-target": int(plan.e_target),
        "smart-utility": float(plan.utility),
        "smart-loss-pre": float(loss_pre),
        "smart-loss-post": float(stats["last_loss"]),
    }
    if env.partial_work:
        # [B] solo con il lavoro parziale acceso, cosi' le altre run mandano
        # esattamente le metriche di prima. smart-weight = campioni visti
        # davvero (n se ha fatto almeno un'epoca): il peso nella media.
        metrics["smart-partial"] = int(plan.steps_cap > 0)
        metrics["smart-weight"] = int(min(n, steps * plan.batch))
    content = RecordDict({"arrays": ArrayRecord(model.state_dict()),
                          "metrics": MetricRecord(metrics)})
    return Message(content=content, reply_to=msg)
\end{lstlisting}

\begin{lstlisting}[style=py,firstnumber=173,caption={\texttt{tesifl/tesiFL/client\_policy.py}, lines 173--226}]
def decide(reading: DeviceReading, env: Envelope, n_samples: int,
           loss: float, n_classes: int = 10) -> Plan:
    """Sceglie (E, B) dentro l'envelope. Vedi la docstring del modulo."""
    u = statistical_utility(loss, n_classes)
    h = headroom(reading, env.soc_min)
    span = max(0, env.e_hi - env.e_lo)
    e_target = env.e_lo + int(round(span * u * h))
    p = env.probe_samples

    def options(e: int):
        out = []
        for b in env.batches:
            t, soc_end = reading.predict(e, b, n_samples, p)
            if _energy_ok(reading, soc_end, env.soc_min):
                out.append((b, t, soc_end))
        return out

    # B dal piu' piccolo: dentro ogni B si riduce E, e si sale di B solo se
    # a quel B non ci sta nemmeno E_lo
    for b in sorted(env.batches):
        for e in range(e_target, env.e_lo - 1, -1):
            t, soc_end = reading.predict(e, b, n_samples, p)
            if _energy_ok(reading, soc_end, env.soc_min) and t <= env.deadline_s:
                return Plan(False, e, b, t, soc_end, True, u, e_target,
                            "target" if e == e_target else "ridotto")

    # [B] 06/11, LAVORO PARZIALE (env.partial_work): nessun E >= E_lo sta
    # nella deadline. Invece di fare E_lo in ritardo (il round aspetta e tutti
    # pagano idle), il client fa i passi di SGD che stanno nella deadline, col
    # B piu' piccolo che la batteria regge (piu' passi per lo stesso tempo,
    # come nel resto della regola). Riguarda soprattutto i telefoni lenti con
    # molti dati, che non stanno nella deadline nemmeno a riposo: col solo
    # rifiuto per deadline non partecipavano mai, accettando arrivavano
    # sempre in ritardo (workload x1: ~430 ritardi su 495). L'update pesa nella
    # media per i campioni visti davvero (smart-weight, vedi sage_smart).
    # Solo se E_lo e' sostenibile per la batteria (feasible): il caso
    # "riserva" (ritiro) resta quello di prima.
    feasible = options(env.e_lo)
    if feasible and env.partial_work:
        for b in sorted(env.batches):
            k = deadline_steps(reading, env, n_samples, b)
            t, soc_end = reading.predict(k / steps_per_epoch(n_samples, b), b, n_samples, p)
            if _energy_ok(reading, soc_end, env.soc_min):
                return Plan(False, env.e_lo, b, t, soc_end, t <= env.deadline_s, u, e_target,
                            "parziale", steps_cap=k)

    # nessun E rispetta la deadline: E_lo col batch piu' veloce, se la
    # batteria lo regge. La deadline e' un vincolo di comparabilita', la
    # riserva no.
    if feasible:
        b, t, soc_end = min(feasible, key=lambda o: o[1])
        return Plan(False, env.e_lo, b, t, soc_end, False, u, e_target,
                    "fuori-deadline")
    return Plan(True, utility=u, e_target=e_target, reason="riserva")
\end{lstlisting}

\begin{lstlisting}[style=py,firstnumber=229,caption={\texttt{tesifl/tesiFL/client\_policy.py}, lines 229--283}]
def steps_per_epoch(n_samples: int, batch: int) -> int:
    """ceil(n / B), come len(DataLoader) nel client."""
    return max(1, math.ceil(n_samples / batch))


def deadline_steps(reading: DeviceReading, env: Envelope, n_samples: int, batch: int) -> int:
    """Massimo numero di passi di SGD (al piu' E_lo epoche) la cui durata,
    stima della loss compresa e col workload attuale, sta nella deadline.
    Almeno 1: se non ci sta nemmeno un passo, se ne fa uno (ritardo minimo).
    La durata e' lineare nelle epoche, quindi la ricerca binaria e' esatta."""
    spe = steps_per_epoch(n_samples, batch)

    def fits(k: int) -> bool:
        return reading.work_time_s(k / spe, batch, n_samples, env.probe_samples) <= env.deadline_s

    hi = env.e_lo * spe
    if fits(hi):
        return hi
    lo = 0                       # fits(0): la sola stima della loss, si assume
    while hi - lo > 1:
        mid = (lo + hi) // 2
        if fits(mid):
            lo = mid
        else:
            hi = mid
    return max(1, lo)


def affordable_steps(reading: DeviceReading, env: Envelope, n_samples: int,
                     batch: int, cap: int) -> int:
    """Massimo numero di passi di SGD che lasciano SoC finale >= soc_min.

    E' il S_max della proposta (budget energetico). Serve da guardia durante
    il training: se la previsione fosse ottimistica, il client si fermerebbe
    qui e invierebbe un update parziale invece di intaccare la riserva.
    """
    if reading.charging:
        return cap
    spe = max(1, math.ceil(n_samples / batch))

    def ok(steps: int) -> bool:
        _, soc_end = reading.predict(steps / spe, batch, n_samples,
                                     env.probe_samples)
        return soc_end >= env.soc_min

    if ok(cap):
        return cap
    lo, hi = 0, cap          # ok(lo) e' vero se il minimo era sostenibile
    while hi - lo > 1:
        mid = (lo + hi) // 2
        if ok(mid):
            lo = mid
        else:
            hi = mid
    return lo
\end{lstlisting}

\begin{lstlisting}[style=py,firstnumber=41,caption={\texttt{tesifl/device/fuel\_gauge.py}, lines 41--130}]
@dataclass(frozen=True)
class DeviceReading:
    """Istantanea del device presa al momento dell'esecuzione (late binding)."""

    soc: float
    workload: float
    charging: bool
    profile: DeviceProfile
    model: str = "nm"

    # -------------------------------------------------- lato mondo (server)
    @classmethod
    def from_world(cls, world: WorldState, cid: int) -> "DeviceReading":
        return cls(
            soc=float(world.snapshot(cid).soc),
            workload=float(world.utilization(cid)),
            charging=bool(world.is_charging(cid)),
            profile=world.profile(cid),
            model=world.battery,
        )

    def to_record(self) -> dict:
        p = self.profile
        return {
            "soc": float(self.soc),
            "workload": float(self.workload),
            "charging": bool(self.charging),
            "battery-model": _NAMES[self.model],
            "cid": int(p.cid),
            "capacity-mah": float(p.battery_capacity_mah),
            "macs-per-s": float(p.macs_per_s),
            "train-power-w": float(p.train_power_w),
            "peukert-n": float(p.peukert_n),
            "tier": str(p.tier),
        }

    # -------------------------------------------------- lato device (client)
    @classmethod
    def from_record(cls, rec) -> "DeviceReading":
        return cls(
            soc=float(rec["soc"]),
            workload=float(rec["workload"]),
            charging=bool(rec["charging"]),
            profile=DeviceProfile(
                cid=int(rec["cid"]),
                battery_capacity_mah=float(rec["capacity-mah"]),
                macs_per_s=float(rec["macs-per-s"]),
                peukert_n=float(rec["peukert-n"]),
                tier=str(rec["tier"]),
                train_power_w=float(rec["train-power-w"]),
            ),
            model={v: k for k, v in _NAMES.items()}[str(rec["battery-model"])],
        )

    # ------------------------------------------------------------ previsioni
    def work_time_s(self, epochs: float, batch_size: int, n_samples: int,
                    probe_samples: int = 0) -> float:
        """Durata del lavoro locale col workload ATTUALE.

        `epochs` puo' essere frazionario: un training interrotto a meta'
        epoca dura la frazione corrispondente (training_time_s e' lineare
        nelle epoche).
        """
        p = effective_profile(self.profile, self.workload)
        return (training_time_s(p, epochs, n_samples, batch_size)
                + probe_time_s(p, probe_samples))

    def soc_after(self, dt_s: float, batch_size: int) -> tuple[float, float]:
        """(SoC dopo il training, SoC dopo l'upload) secondo la batteria del
        device. In carica la batteria non scende: e' collegato alla rete.
        La potenza di training e' quella del tier: `batch_size` cambia solo la
        durata (gia' dentro dt_s), non la potenza."""
        if self.charging:
            return self.soc, self.soc
        s1 = self._drain(self.soc, training_power_w(self.profile), dt_s)
        if s1 <= 0.0:
            # il mondo marca il device come morto e non addebita l'upload
            return 0.0, 0.0
        s2 = self._drain(s1, comm_power_w(), comm_time_s())
        return s1, s2

    def predict(self, epochs: float, batch_size: int, n_samples: int,
                probe_samples: int = 0) -> tuple[float, float]:
        """(durata, SoC finale dopo l'upload) di un piano (E, B)."""
        dt = self.work_time_s(epochs, batch_size, n_samples, probe_samples)
        return dt, self.soc_after(dt, batch_size)[1]

    def _drain(self, soc: float, power_w: float, dt_s: float) -> float:
        d = battery_delta_soc(self.model, power_w, dt_s, self.profile, soc)
        return max(0.0, min(1.0, soc + d))
\end{lstlisting}

\begin{lstlisting}[style=py,firstnumber=5,caption={\texttt{tesifl/tesiFL/task.py}, lines 5--80}]
def train(net, trainloader, epochs, lr, device, proximal_mu: float = 0.0,
          max_steps: int | None = None, stats: dict | None = None):
    """Train the model on the training set.

    proximal_mu > 0 attiva il termine prossimale di FedProx:
        loss = CE + (mu/2) * ||w - w_global||^2
    dove w_global sono i pesi ricevuti dal server a inizio round (congelati,
    detached: il gradiente non ci passa attraverso). mu=0 -> FedAvg puro,
    nessun costo computazionale aggiuntivo.

    max_steps  interrompe il training dopo quel numero di passi di SGD, anche
               a meta' epoca (Fase 6 della proposta: early stopping quando il
               budget energetico finisce). None = epoche complete, come prima.
    stats      se e' un dict, ci scrive "steps" (passi eseguiti) e
               "last_loss" (loss media dell'ultima epoca, anche parziale):
               servono a sage_smart. Il valore di ritorno non cambia.
    """
    net.to(device)
    criterion = torch.nn.CrossEntropyLoss().to(device)
    optimizer = torch.optim.SGD(net.parameters(), lr=lr, momentum=0.9)

    # [B] snapshot dei pesi globali PRIMA di qualunque step locale
    use_prox = proximal_mu > 0.0
    if use_prox:
        global_params = [p.detach().clone() for p in net.parameters()]

    # [B] reduction='none' per avere la loss dei singoli campioni: serve alla
    # statistical utility di Oort (|B|*sqrt(mean(Loss^2))). Il backward usa
    # comunque la media, quindi FedAvg/FedProx sono numericamente invariati.
    criterion_ns = torch.nn.CrossEntropyLoss(reduction="none").to(device)

    net.train()
    running_loss = 0.0
    sq_sum, n_seen = 0.0, 0
    steps = 0
    last_sum, last_n = 0.0, 0
    for _ep in range(epochs):
        # [B] controllo PRIMA dell'azzeramento: se il budget finisce
        # esattamente a fine epoca, le statistiche restano quelle dell'epoca
        # appena chiusa invece di un'epoca vuota
        if max_steps is not None and steps >= max_steps:
            break
        # statistiche riferite all'ULTIMA epoca: riflettono lo stato del
        # modello a fine round, non la media del percorso
        sq_sum, n_seen = 0.0, 0
        last_sum, last_n = 0.0, 0
        for batch in trainloader:
            if max_steps is not None and steps >= max_steps:
                break
            images = batch["img"].to(device)
            labels = batch["label"].to(device)
            optimizer.zero_grad()
            per_sample = criterion_ns(net(images), labels)
            with torch.no_grad():
                sq_sum += float(torch.sum(per_sample ** 2).item())
                n_seen += int(per_sample.numel())
                last_sum += float(torch.sum(per_sample).item())
                last_n += int(per_sample.numel())
            loss = per_sample.mean()
            if use_prox:
                prox_term = 0.0
                for p, g in zip(net.parameters(), global_params):
                    prox_term = prox_term + torch.sum((p - g) ** 2)
                loss = loss + (proximal_mu / 2.0) * prox_term
            loss.backward()
            optimizer.step()
            running_loss += loss.item()
            steps += 1
    # [B] diviso per i passi EFFETTIVI: con epoche complete steps vale
    # epochs * len(trainloader), cioe' esattamente il denominatore di prima
    avg_trainloss = running_loss / max(1, steps)
    if stats is not None:
        stats["steps"] = steps
        stats["last_loss"] = last_sum / last_n if last_n else float("nan")
    # (loss media, somma dei quadrati per campione, n campioni dell'ultima epoca)
    return avg_trainloss, sq_sum, n_seen
\end{lstlisting}

\newpage
%======================================================================
\section{Step 7: accounting, state update and aggregation}
\label{sec:aggregate}

For each reply, in order of phone index:
\begin{itemize}
  \item Work done (\cf{\_work\_done}): duration $t_i = t^{\mathrm{tr}}_i(\nu_i /
        \lceil n_i/B \rceil, n_i, B) + t^{\mathrm{probe}}_i(n_{\mathrm{p}})$ with
        the workload of this round, from the steps that the phone declares.
  \item Training energy: power $P^{\mathrm{tr}}_i$ for $t_i$ on the battery of
        the world (\cf{apply\_round}); nothing on the battery if charging.
  \item Communication: model download and upload ($P^{\mathrm{comm}}$ for
        $t^{\mathrm{comm}}$); download and an isolated control message for a
        phone that withdrew. The communication time does not enter the round
        duration.
  \item State update (\cf{\_on\_round\_result}): the class histogram updates
        $D$; the reported final SoC, workload and charging flag become the last
        report, at time $t^{\mathrm{now}} + t^{\mathrm{ovh}}_r + t_i$; $n_i$ is
        updated; if the phone was not charging and did at least one step, the
        calibration factor is updated with the observed drop:
        \begin{equation}
          k_i \leftarrow (1 - \alpha)\, k_i + \alpha \, \frac{s^{\mathrm{start}}_i - s^{\mathrm{end}}_i}
          {\bigl(P^{\mathrm{tr}}_i\, t'_i + P^{\mathrm{comm}} t^{\mathrm{comm}}\bigr) / (3600\, \eta\, V_0 Q_i)} ,
        \end{equation}
        with $t'_i$ the training and probe time recomputed with the declared
        workload, applied only if numerator and denominator are positive.
\end{itemize}
Then the round duration $t_r = \max_{i \in S_r} t_i + t^{\mathrm{ovh}}_r$ is
fixed, the server clock advances by $t_r$, the charging phones recharge, and
every alive phone pays the idle for $t_r - (t_i + t^{\mathrm{ovh}}_r)$ if
selected and $t_r$ otherwise.

The update of a phone that withdrew, or whose $(E, B)$ is outside its envelope
($E \notin [E_{\mathrm{lo}}, E_{\mathrm{hi},i}]$ or $B \notin \mathcal{B}$), is
not averaged; its energy is charged anyway. The other updates are averaged with
the Flower \cf{FedAvg}, weighted by \cf{num-examples} $= n_i$. If at least two
updates are valid and in this round at least one phone did partial work
($\min(n_i, \nu_i B) < n_i$), the average is replaced by the one weighted by
the samples really seen:
\begin{equation}
  x^{r+1} = \sum_{i} q_i \, x^{r}_i , \qquad
  q_i = \frac{\min(n_i, \nu_i B_i)}{\sum_j \min(n_j, \nu_j B_j)} ,
\end{equation}
in double precision, BatchNorm buffers included; integer buffers are rounded.
A phone that did at least one full epoch has $\min(n_i, \nu_i B_i) = n_i$, so
the two averages differ only in the weight of the partial updates. The columns
\cf{tau\_min}, \cf{tau\_max}, \cf{tau\_eff} and \cf{tau\_spread} record the
declared steps and do not change the average.

\begin{lstlisting}[style=py,firstnumber=777,caption={\texttt{tesifl/strategies/sage\_smart.py}, lines 777--842}]
    def _rep(self, reply: Message, key: str, default=None):
        return self._metric(reply, key, default=default)

    def _is_optout(self, reply: Message) -> bool:
        return int(self._rep(reply, "smart-optout", 0) or 0) == 1

    def _charge_communication(self, cid: int, reply: Message) -> tuple[float, float]:
        """Chi si ritira ha scaricato il modello ma non carica l'update: paga
        il download e la risposta di ritiro (uno scambio di controllo
        isolato), non lo scambio intero. [B] 04/11: prima pagava ~181 J invece
        di ~48 J, e l'ablazione senza negoziazione risultava penalizzata."""
        if not self._is_optout(reply):
            return super()._charge_communication(cid, reply)
        e_dl, l_dl = self.world.apply_download(cid)
        e_ct, l_ct = self.world.apply_control(cid, merged=False)
        return e_dl + e_ct, l_dl + l_ct

    def _work_done(self, cid: int, n_ex: int, reply: Message) -> tuple[float, int]:
        """Lavoro dichiarato dal client -> durata, per l'addebito fisico.

        Stessa formula che il client usa per prevedere il proprio SoC
        (DeviceReading.work_time_s).
        """
        b = int(self._rep(reply, "smart-batch", self._batch_for(cid)))
        steps = int(self._rep(reply, "smart-steps", 0))
        probe = int(self._rep(reply, "smart-probe", 0))
        e_eff = steps / math.ceil(n_ex / b) if steps > 0 and n_ex > 0 else 0.0
        dt = (self.world.round_duration_s(cid, epochs=e_eff, batch_size=b,
                                          dataset_size_local=n_ex)
              + self.world.probe_duration_s(cid, probe))
        return dt, b

    def _in_envelope(self, cid: int, reply: Message) -> bool:
        env = self._plan.get(cid)
        if env is None:
            return False
        e = int(self._rep(reply, "smart-epochs", 0))
        b = int(self._rep(reply, "smart-batch", 0))
        return env.e_lo <= e <= env.e_hi and b in env.batches

    def _aggregatable(self, replies: list[Message]) -> list[Message]:
        """Esclude chi si e' ritirato e chi e' uscito dall'envelope.

        [B] il client non dovrebbe mai uscire dall'envelope: la sua policy
        sceglie dentro i limiti. Il controllo rende esplicito il contratto:
        fuori dai limiti il server non accetta l'update (ma il lavoro e'
        stato fatto, quindi l'energia viene addebitata lo stesso).
        """
        out = []
        for r in replies:
            if not r.has_content():
                out.append(r)      # gli errori li gestisce FedAvg, come prima
                continue
            if self._is_optout(r):
                continue
            cid = self._node_to_cid.get(int(r.metadata.src_node_id))
            if cid is not None and not self._in_envelope(cid, r):
                continue
            out.append(r)
        return out

    def _steps(self, reply: Message) -> float:
        """Passi di SGD dichiarati dal client (colonne tau_* del CSV)."""
        if self._is_optout(reply):
            return 0.0
        return float(self._rep(reply, "smart-steps", 0))
\end{lstlisting}

\begin{lstlisting}[style=py,firstnumber=844,caption={\texttt{tesifl/strategies/sage\_smart.py}, lines 844--906}]
    def _on_round_result(
        self, cid: int, n_ex: int, dt_s: float, e_wh: float, e_lin_wh: float,
        reply: Message, server_round: int,
    ) -> None:
        super()._on_round_result(cid, n_ex, dt_s, e_wh, e_lin_wh, reply,
                                 server_round)
        optout = self._is_optout(reply)
        soc0 = float(self._rep(reply, "smart-soc-start", self._resource(cid)))
        soc1 = float(self._rep(reply, "smart-soc-final", soc0))
        chg = bool(int(self._rep(reply, "smart-charging", 0)))
        w = float(self._rep(reply, "smart-workload", self.w_busy))
        b = int(self._rep(reply, "smart-batch", self._batch_for(cid)))
        steps = int(self._rep(reply, "smart-steps", 0))
        probe = int(self._rep(reply, "smart-probe", 0))

        # report del client, con l'istante in cui e' stato fatto: fine del
        # suo training e upload, dopo il ritardo iniziale del round. Da qui
        # partono le previsioni dei round successivi (passo 1).
        self._record(cid, soc1, w, chg, self._clock_s + self._overhead_s + dt_s)
        if n_ex > 0:
            self._n[cid] = int(n_ex)
        loss = self._rep(reply, "smart-loss-pre", None)
        if loss is not None:
            self._loss_known[cid] = float(loss)

        # calibrazione k_i: calo di SoC dichiarato / stima lineare dello
        # stesso lavoro, con la durata ricostruita dal workload DICHIARATO.
        # [B] nel mondo lineare k_i resta ~1; con Peukert e datasheet impara
        # quanto la batteria del device paga il training rispetto a P*t.
        if not chg and steps > 0 and n_ex > 0:
            prof = self.world.profile(cid)
            peff = effective_profile(prof, w)
            e_eff = steps / math.ceil(n_ex / b)
            t = training_time_s(peff, e_eff, n_ex, b) + probe_time_s(peff, probe)
            cap_wh = prof.battery_capacity_mah / 1000.0 * V_NOMINAL
            lin = (energy_wh_linear(training_power_w(prof), t)
                   + energy_wh_linear(comm_power_w(), comm_time_s())) / cap_wh
            drop = soc0 - soc1
            if lin > 0 and drop > 0:
                a = self.calib_alpha
                self._k_cal[cid] = (1 - a) * self._k_cal.get(cid, 1.0) + a * drop / lin

        self._reports[cid] = {
            "optout": optout,
            "epochs": int(self._rep(reply, "smart-epochs", 0)),
            "e_eff": (steps / math.ceil(n_ex / b)) if steps > 0 and n_ex > 0 else 0.0,
            "batch": b,
            "steps": steps,
            "budget_hit": int(self._rep(reply, "smart-budget-hit", 0)),
            "deadline_ok": int(self._rep(reply, "smart-deadline-ok", 1)),
            "charging": chg,
            "loss_pre": float(loss) if loss is not None else float("nan"),
            "loss_post": float(self._rep(reply, "smart-loss-post", float("nan"))),
            "workload": w,
            # [B] VERIFICA, non input: SoC dichiarato contro SoC VERO del
            # mondo dopo l'addebito di training + upload (l'idle viene dopo)
            "soc_err": abs(float(self.world.snapshot(cid).soc) - soc1),
            "rejected": (not optout) and not self._in_envelope(cid, reply),
            "partial": int(self._rep(reply, "smart-partial", 0)),
        }
        if self._reports[cid]["rejected"]:
            log(INFO, "%s round %s: update del client %s fuori envelope, "
                "scartato", self.variant, server_round, cid)
\end{lstlisting}

\begin{lstlisting}[style=py,firstnumber=965,caption={\texttt{tesifl/strategies/sage\_smart.py}, lines 965--1056}]
    def aggregate_train(self, server_round, replies):
        """Media FedAvg pesata per i dati; col lavoro parziale, per i campioni
        visti davvero.

        La media per i dati la fa super() (PhysicalFedAvg) sulle reply di
        _aggregatable (senza opt-out e fuori envelope), BatchNorm compresa.
        Qui si aggiungono:
          - le colonne della policy nel CSV (_policy_stats);
          - le colonne tau_* coi passi di SGD dichiarati dai client
            (tau_min, tau_max, tau_eff = media pesata per i dati, tau_spread
            = max/min): solo diagnostica, nessun effetto sulla media;
          - [B] 06/11 LAVORO PARZIALE: nei round in cui qualcuno ha fatto meno
            di un'epoca, la media e' pesata per i campioni visti davvero,
            min(n, passi * B), invece che per n: un update da meta' dati non
            deve pesare come uno da tutti. Negli altri round resta la media di
            super(), identica.
        [B] 16/11: FedNova (normalizzazione per i passi, Wang et al. NeurIPS
        2020) tolto. Era spento dal 30/10: contributo esterno, e nelle due
        ablazioni la media pesata per i dati andava uguale o meglio (FedNova
        -0.9 / -1.0 punti).
        """
        arrays, metrics = super().aggregate_train(server_round, replies)
        # il server sa quanto e' durato il round: ha aspettato le risposte
        self._clock_s += self._last_round_s

        # [B] metriche della policy nel CSV per round, PRIMA dell'uscita
        # anticipata: servono anche quando l'aggregazione non viene applicata.
        # super() ha gia' assegnato _round_stats[round], quindi qui si aggiunge
        # sopra con update() invece di sovrascrivere.
        if self._plan:
            self._round_stats.setdefault(int(server_round), {}).update(
                self._policy_stats())

        if arrays is None or self._global_arrays is None:
            return arrays, metrics

        # [B] le CHIAVI vanno conservate: costruire un ArrayRecord da una lista
        # gli fa assegnare nomi posizionali "0","1",... e il load_state_dict in
        # global_evaluate non li riconosce. Leggere per chiave garantisce anche
        # lo stesso ordine fra record globale e locali.
        keys = list(self._global_arrays.keys())
        rows = []
        wts = []        # campioni visti davvero (smart-weight), per il lavoro parziale
        # stesse reply che super() ha mediato: senza opt-out e fuori envelope
        for r in self._aggregatable(replies):
            if not r.has_content():
                continue
            if self._node_to_cid.get(int(r.metadata.src_node_id)) is None:
                continue
            n = float(self._metric(r, "num-examples", "num_examples", default=0))
            tau = self._steps(r)
            if n <= 0 or tau <= 0:
                continue
            key = sorted(r.content.array_records)[0]
            rec = r.content.array_records[key]
            rows.append((n, tau, [rec[k].numpy() for k in keys]))
            wts.append(float(self._metric(r, "smart-weight", default=n)))

        if len(rows) < 2:
            return arrays, metrics

        tot = sum(n for n, _, _ in rows)
        p = [n / tot for n, _, _ in rows]
        tau_eff = sum(pi * tau for pi, (_, tau, _) in zip(p, rows))
        taus = [tau for _, tau, _ in rows]
        self._round_stats.setdefault(int(server_round), {}).update({
            "tau_min": int(min(taus)),
            "tau_max": int(max(taus)),
            "tau_eff": float(tau_eff),
            "tau_spread": float(max(taus) / max(1.0, min(taus))),
        })
        log(INFO, "round %s: FedAvg su %s client | tau %s-%s, tau_eff %.1f, "
            "spread %.1f", server_round, len(rows), int(min(taus)),
            int(max(taus)), tau_eff, max(taus) / max(1.0, min(taus)))

        if self.partial_work and any(w < n for w, (n, _, _) in zip(wts, rows)):
            tot_w = sum(wts)
            q = [w / tot_w for w in wts]
            merged = ArrayRecord()
            for j, k in enumerate(keys):
                ref = self._global_arrays[k].numpy()
                avg = np.asarray(sum(qi * row[2][j].astype(np.float64) for qi, row in zip(q, rows)),
                                 dtype=np.float64).reshape(ref.shape)
                if np.issubdtype(ref.dtype, np.integer):
                    avg = np.rint(avg)
                merged[k] = Array(avg.astype(ref.dtype))
            self._round_stats.setdefault(int(server_round), {}).update({
                "smart_weight_frac": float(tot_w / sum(n for n, _, _ in rows)),
            })
            return merged, metrics

        return arrays, metrics
\end{lstlisting}

The accuracy of the CSV is the one of the global model on the test set,
computed by the server after every round. The final accuracy of the reports is
the mean of the last 20 rounds.

\begin{lstlisting}[style=py,firstnumber=421,caption={\texttt{tesifl/tesiFL/server\_app.py}, lines 421--439}]
def global_evaluate(server_round: int, arrays: ArrayRecord) -> MetricRecord:
    """Evaluate model on central data."""

    # Load the model and initialize it with the received weights
    model = build_model(_MODEL_NAME, num_classes(_DATASET))
    model.load_state_dict(arrays.to_torch_state_dict())
    device = torch.device("cuda:0" if torch.cuda.is_available() else "cpu")
    model.to(device)

    # Load entire test set
    test_dataloader = load_centralized_dataset(_DATASET)

    # Evaluate the global model on the test set
    test_loss, test_acc = test(model, test_dataloader, device)

    _ACC_BY_ROUND[int(server_round)] = (float(test_acc), float(test_loss))

    # Return the evaluation metrics
    return MetricRecord({"accuracy": test_acc, "loss": test_loss})
\end{lstlisting}

\newpage
%======================================================================
\section{Logged metrics}
\label{sec:metrics}

Besides the state of the world (energy by item, SoC distribution, depleted and
charging phones, \cf{round\_s}, \cf{sim\_time\_s}), every round adds:
\begin{longtable}{p{0.30\textwidth}p{0.64\textwidth}}
\toprule
Column & Content \\
\midrule
\endhead
\cf{mean\_epochs}, \cf{max\_epochs}, \cf{smart\_min\_epochs} & epochs really executed, $\nu_i / \lceil n_i/B \rceil$, over the phones that did not withdraw \\
\cf{mean\_batch} & mean batch size chosen \\
\cf{n\_charging\_sel} & selected phones that were charging \\
\cf{smart\_mean\_e\_hi} & mean $E_{\mathrm{hi}}$ of the envelopes of the selected phones \\
\cf{smart\_n\_bound\_energy}, \cf{smart\_n\_bound\_deadline} & envelopes limited by the energy or by the deadline, counted over all the envelopes computed in the round (proposals included) \\
\cf{smart\_n\_optout}, \cf{smart\_n\_rejected} & withdrawals, updates outside the envelope \\
\cf{smart\_n\_deadline\_miss}, \cf{smart\_n\_budget\_hit} & trainings beyond the deadline, trainings stopped by the energy budget \\
\cf{smart\_mean\_loss\_pre}, \cf{smart\_mean\_loss\_post} & probe loss, loss of the last epoch \\
\cf{smart\_mean\_workload} & mean reported workload of the selected phones \\
\cf{smart\_soc\_err} & maximum $|$true SoC $-$ reported final SoC$|$ (check) \\
\cf{smart\_mean\_k} & mean $k_i$ \\
\cf{smart\_deadline\_s} & $T$ \\
\cf{smart\_info\_err} & mean $|\hat{s}_i - s_i|$ over $\mathcal{A}_r$ (check) \\
\cf{smart\_n\_queried}, \cf{smart\_n\_refused}, \cf{smart\_waves} & proposals that reached a phone, refusals, waves \\
\cf{smart\_overhead\_s} & $t^{\mathrm{ovh}}_r$ \\
\cf{smart\_n\_partial} & phones that did partial work \\
\cf{smart\_n\_deferred}, \cf{smart\_n\_noreply}, \cf{smart\_w\_err} & deferred phones, unreachable phones, mean $|\hat{w}_i - w_i|$ (check) \\
\cf{tau\_*}, \cf{smart\_weight\_frac} & declared steps; fraction of samples seen in the rounds with partial work \\
\bottomrule
\end{longtable}

\begin{lstlisting}[style=py,firstnumber=908,caption={\texttt{tesifl/strategies/sage\_smart.py}, lines 908--962}]
    def _policy_stats(self) -> dict:
        """Colonne della policy nel CSV per round."""
        rep = list(self._reports.values())
        done = [r for r in rep if not r["optout"]]
        envs = list(self._plan.values())

        def mean(xs):
            xs = [x for x in xs if x == x]          # scarta i NaN
            return float(np.mean(xs)) if xs else float("nan")

        st = {
            # epoche e batch ESEGUITI dai client
            "mean_epochs": mean([r["e_eff"] for r in done]),
            "max_epochs": float(max((r["e_eff"] for r in done), default=0.0)),
            "mean_batch": mean([r["batch"] for r in done]),
            "n_charging_sel": int(sum(1 for r in rep if r["charging"])),
            "smart_min_epochs": float(min((r["e_eff"] for r in done), default=0.0)),
            "smart_mean_e_hi": mean([e.e_hi for e in envs]),
            "smart_n_bound_energy": int(sum(1 for v in self._bound.values() if v == "energy")),
            "smart_n_bound_deadline": int(sum(1 for v in self._bound.values() if v == "deadline")),
            "smart_n_optout": int(sum(1 for r in rep if r["optout"])),
            "smart_n_rejected": int(sum(1 for r in rep if r["rejected"])),
            "smart_n_deadline_miss": int(sum(1 for r in done if not r["deadline_ok"])),
            "smart_n_budget_hit": int(sum(1 for r in done if r["budget_hit"])),
            "smart_mean_loss_pre": mean([r["loss_pre"] for r in rep]),
            "smart_mean_loss_post": mean([r["loss_post"] for r in done]),
            "smart_mean_workload": mean([r["workload"] for r in rep]),
            "smart_soc_err": float(max((r["soc_err"] for r in rep), default=0.0)),
            "smart_mean_k": mean(list(self._k_cal.values())),
            "smart_deadline_s": float(self._deadline_s),
            # passi 1 e 4: qualita' e costo dello stato usato per selezionare
            "smart_info_err": float(self._info_err),
            "smart_n_queried": int(self._n_queried),
            "smart_n_refused": int(self._n_refused),
            "smart_waves": int(self._waves),
            "smart_overhead_s": float(self._overhead_s),
        }
        if self.partial_work:
            # [B] solo con il lavoro parziale (nopartial ha le colonne di prima)
            st["smart_n_partial"] = int(sum(1 for r in done if r.get("partial")))
        if self._protocol_cols:
            # [B] rinviati (accettato ma non scelto), non raggiunti senza
            # broadcast, errore medio del workload noto (verifica)
            st["smart_n_deferred"] = int(self._n_deferred)
            st["smart_n_noreply"] = int(self._n_noreply)
            st["smart_w_err"] = float(self._w_err)
        log(INFO, "%s: %s risposte | E eseguite %.1f (min %.2f, max %.2f) "
            "| E_hi medio %.1f (energia %s, deadline %s) | opt-out %s | "
            "fuori deadline %s | errore SoC %.1e | errore stato noto %.3f",
            self.variant, len(rep), st["mean_epochs"], st["smart_min_epochs"],
            st["max_epochs"], st["smart_mean_e_hi"], st["smart_n_bound_energy"],
            st["smart_n_bound_deadline"], st["smart_n_optout"],
            st["smart_n_deadline_miss"], st["smart_soc_err"],
            st["smart_info_err"])
        return st
\end{lstlisting}

Observed values (datasheet battery, seed 2026, rounds 21--250 for the protocol
columns):

\begin{table}[H]
  \centering
  \begin{tabular}{lccc}
    \toprule
    & $\beta = 0.1$ & $\beta = 0.5$ & $\beta = 1.0$ \\
    \midrule
    Deadline $T$ (s)                          & 132.4 & 130.5 & 127.9 \\
    Proposals that reached a phone, per round & 26.8  & 26.9  & 27.0 \\
    Unreachable phones per round              & 2.90  & 2.92  & 2.93 \\
    Refusals per round                        & 0.21  & 0.18  & 0.07 \\
    Deferred phones per round                 & 6.89  & 6.89  & 7.00 \\
    Waves per round                           & 1.00  & 1.00  & 1.00 \\
    Control overhead per round (s)            & 0.97  & 0.97  & 0.97 \\
    Epochs executed (mean, all rounds)        & 2.51  & 2.25  & 2.21 \\
    Batch size (mean, all rounds)             & 16.1  & 16.1  & 17.5 \\
    Partial trainings per round (all rounds)  & 2.44  & 2.08  & 1.58 \\
    Withdrawals in the run                    & 0     & 0     & 0 \\
    \bottomrule
  \end{tabular}
\end{table}

\newpage
%======================================================================
\section{Ablation switches and labels}
\label{sec:ablation}

The run parameter \cf{smart-ablation} switches components off, several of them
joined by \cf{+}. The tokens are the keys of the dictionary in
\cf{server\_app.py} (Section~\ref{sec:params}).

\begin{longtable}{p{0.13\textwidth}p{0.81\textwidth}}
\toprule
Token & Effect \\
\midrule
\endhead
\cf{nostale} & $c = 0$: no staleness, in the ranking and in the fresh score \\
\cf{noest} & no forecast: $\hat{s}_i$ and $\hat{w}_i$ are the last report as it is \\
\cf{noneg} & no proposals and no fresh choice: the model goes to the first $h_r$ phones of the ranking, which can still withdraw after the download. The broadcast of the state comes back (\cf{broadcast = true} in \cf{experiment.toml}): every alive phone answers with probability 0.9, the others are excluded from the round, the answers pay a control message (shared with the download for the selected phones), and the round waits $t_{\mathrm{ctrl}}$, or $\max(t_{\mathrm{ctrl}}, t_{\mathrm{to}}) = 1$ s if someone does not answer \\
\cf{e5} & $E_{\mathrm{lo}} = E_{\mathrm{cap}} = 5$: five epochs fixed \\
\cf{b128} & $\mathcal{B} = \{128\}$ \\
\cf{nopartial} & no partial work: a phone that does not fit the deadline with $E_{\mathrm{lo}}$ trains $E_{\mathrm{lo}}$ epochs out of the deadline \\
\bottomrule
\end{longtable}

The labels of the cumulative ablations keep the names of the runs already done:

\begin{table}[H]
  \centering
  \begin{tabular}{llp{0.42\textwidth}}
    \toprule
    Label & Tokens & Step \\
    \midrule
    \cf{sage\_smart\_cum0}   & \cf{nostale+noest+noneg+e5+b128+nopartial} & 0 base (both orders) \\
    \cf{sage\_smart\_cum1}   & \cf{noest+noneg+e5+b128+nopartial} & 1 + staleness (both orders) \\
    \cf{sage\_smart\_cum2p}  & \cf{noest+noneg+b128} & forward 2: + adaptive $E$ with partial work \\
    \cf{sage\_smart\_pcum4}  & \cf{noest+noneg} & forward 3: + adaptive $B$ \\
    \cf{sage\_smart2\_nofc}  & \cf{noest} & forward 4: + negotiation \\
    \cf{sage\_smart2\_rev2}  & \cf{e5+b128+nopartial} & reverse 2: + negotiation and forecast \\
    \cf{sage\_smart2\_rev3}  & \cf{b128} & reverse 3: + adaptive $E$ \\
    \cf{sage\_smart}         & none & last step of both orders \\
    \bottomrule
  \end{tabular}
\end{table}

\newpage
%======================================================================
\section{What the server knows, and simplifications}
\label{sec:info}

The decisions of the strategy use only:
\begin{itemize}
  \item static data of the enrollment: tier, throughput, training power,
        capacity and battery model of every phone, $n_i$;
  \item population data: $\bar{w}$ (mean of the stationary workloads of the
        three chains), the plug-in rate $\lambda$ and the charging speed
        $\gamma$;
  \item what the phones send: the answers to the proposals and the reports
        after the training;
  \item the durations of the rounds, which the server observes.
\end{itemize}
The world enters only in two places, both outside the decisions: the sensor
record that the server attaches to the messages on behalf of the phone (a real
phone reads it from its operating system), and the check columns
\cf{smart\_soc\_err}, \cf{smart\_info\_err} and \cf{smart\_w\_err}.

Simplifications of the simulator, as implemented:
\begin{itemize}
  \item The forecast uses the same battery and idle models as the world, so it
        is wrong only because of recharges and workload. On real phones it would
        also be wrong because of the model: here it is optimistic.
  \item The enrollment has no energy or time cost; $n_i$ is computed by the
        server with the same partitioner as the clients.
  \item The proposals see the state before \cf{begin\_round}; the training sees
        the state after it (one workload transition and the recharge draws in
        between).
  \item The communication time (71 s per exchange) is charged as energy but
        does not enter the round duration, which is the slowest training plus
        the control overhead.
  \item The time model charges the batch overhead on $\lfloor n/B \rfloor$
        batches per epoch, while the client executes $\lceil n/B \rceil$ steps
        per epoch.
  \item The acceptance check uses 256 probe samples even when $n_i < 256$; the
        training uses $\min(256, n_i)$.
  \item Within a round the power of each item (training, communication, idle)
        is constant.
\end{itemize}

\end{document}
